\section{训练路线章正文案例推导补强}

这一节补齐本章正文出现过的 LeetCode 案例推导。阅读顺序统一按“题目说明、样例入口、暴力基线、暴力过程、重复工作、优化条件、相邻状态、状态复用、指针和字段、代码落点、正确性、边界实验、复杂度变化、迁移追问”展开；目的不是新增题量，而是把正文中的案例都讲成可复述、可证明、可迁移的解题链。

\subsection{LeetCode 1 两数之和}

\textbf{题目说明。} LeetCode 1 两数之和 的题面入口要先还原成具体任务：当前元素只需要寻找唯一补数，历史值到下标的映射就是最小状态。

\textbf{样例入口。} 拿数组 [2, 7]、target=9 手算：站在 2 时历史为空，不能配对；站在 7 时只需要问历史里有没有 2。再试 [3, 3]、target=6：如果先把当前 3 插入再查，就会把同一个下标用两次。

\textbf{暴力基线。} 暴力两层循环枚举 i 和 j，检查 nums[i]+nums[j] 是否等于 target，找到后返回两个下标。

\textbf{暴力过程。} 以 nums=[2, 7, 11, 15]、target=9 为例，暴力先试 2+7 命中；换成 [3, 2, 4]、target=6 时，会试 3+2 失败，再试 2+4 命中。

\textbf{重复工作。} 题面模型：当前元素只需要寻找唯一补数，历史值到下标的映射就是最小状态。暴力版的慢点要落到本题：浪费在每个当前位置都回头扫描历史。站在 nums[i] 时，真正需要的不是所有历史值，而是历史里有没有 target-nums[i]。后面的优化只消掉这类重复工作，不能改变答案集合。

\textbf{优化条件。} 本题的关键观察是：先查再插入，避免同一个下标被用两次；若数组有序才考虑双指针。这句话必须能翻译成可维护状态；如果题目条件改变后观察不再成立，就不能继续套原来的写法。

\textbf{相邻状态。} 规律不是“用哈希表”四个字，而是每个当前位置只和它左侧历史配对，代码必须先查补数再插入当前值。把这个特例继续拆开看：每个合法答案都要还原成当前状态和某个历史状态的配对；哈希表的 key 负责定位历史状态，value 负责表达次数、最早位置或存在性。先查还是先写，必须由是否允许当前前缀和自己配对来决定。

\textbf{状态复用。} 哈希表保存已经扫描过的值到下标的映射。每到一个新数，先查补数是否存在；不存在再把当前数放入历史。手算时只改被新输入影响的那一小部分，而不是回头重算完整候选。

\textbf{指针和字段。} 状态字段要能说清：i 是当前下标，need=target-nums[i] 是补数，seen[value] 是该值在左侧历史中出现的下标。手算时按这个协议跟踪：nums=[3, 3]、target=6 时，第一个 3 先查不到补数，再插入；第二个 3 查到历史 3，返回两个不同下标。

\textbf{代码落点。} 关键顺序是先 find(need)，再 seen[nums[i]]=i。若先插入当前值，target=2*nums[i] 时可能把同一位置用两次。关键代码行必须能逐句翻译成“旧状态怎样变成新状态”，否则就只是模板。

\textbf{正确性。} 每个答案对都有一个较右下标；扫描到较右下标时，较左下标已经在 seen 中，所以一定能被查到。

\textbf{边界实验。} 先用重复数字、目标为偶数且补数等于当前值、没有答案时返回空结果做反例检查。实验时把重复数字、目标为偶数且补数等于当前值、没有答案时返回空结果加入对拍集合，用平方暴力和优化版比较。失败时先看初始历史状态、查询更新顺序和哈希表 value 类型。

\textbf{复杂度变化。} 暴力 O(\ensuremath{n^{2}})，哈希表扫描 O(n) 均摊时间、O(n) 空间。

\textbf{迁移追问。} 如果题目改成“若改成返回所有不重复数对，要先排序或用频次表处理重复”，先判断原状态是否仍然覆盖新答案集合，再决定是微调边界、改 value 语义，还是换算法。

\subsection{LeetCode 283 移动零}

\textbf{题目说明。} LeetCode 283 移动零给定数组 nums，要求原地把所有 0 移到末尾，同时保持非零元素的相对顺序。

\textbf{样例入口。} 样例 nums=[0, 1, 0, 3, 12]，处理后为 [1, 3, 12, 0, 0]；非零元素 1、3、12 的相对顺序不变。

\textbf{暴力基线。} 暴力可以新建数组，先放所有非零元素，再补同样数量的零，最后复制回原数组。

\textbf{暴力过程。} 以 [0, 1, 0, 3, 12] 为例，辅助数组先收集 [1, 3, 12]，再补 [0, 0]。

\textbf{重复工作。} 题面模型：给定数组，原地把所有 0 移到末尾，同时保持非零元素的相对顺序。暴力版的慢点要落到本题：浪费在辅助数组。非零元素的相对顺序由读取顺序天然保证，只需要一个原地写入位置。后面的优化只消掉这类重复工作，不能改变答案集合。

\textbf{优化条件。} 本题的关键观察是：稳定保留非零元素顺序，把零集中到尾部；慢指针写入下一个非零位置，最后补零或用交换。这句话必须能翻译成可维护状态；如果题目条件改变后观察不再成立，就不能继续套原来的写法。

\textbf{相邻状态。} 规律是先稳定压缩非零前缀，再填零；不能在扫描时盲目交换导致顺序变化。把这个特例继续拆开看：先标出已经确认的前缀或后缀，再标出未知区；能被丢掉的候选，必须是以后再也不会改变已确认区域语义的候选。这样才算从例子推出数组不变量，而不是只记一次操作。

\textbf{状态复用。} write 指向下一个非零元素应写入的位置。第一遍把非零元素稳定写到前面，第二遍把 write 之后填 0。手算时只改被新输入影响的那一小部分，而不是回头重算完整候选。

\textbf{指针和字段。} 状态字段要能说清：read 是扫描位置，write 左侧是不含 0 的稳定前缀，write 到末尾最后会被置零。手算时按这个协议跟踪：读到 1 写到 nums[0]，读到 3 写到 nums[1]，读到 12 写到 nums[2]，最后 nums[3]、nums[4] 置 0。

\textbf{代码落点。} 关键代码是 if (nums[read]!=0) nums[write++]=nums[read]; 后面 while(write<n) nums[write++]=0。关键代码行必须能逐句翻译成“旧状态怎样变成新状态”，否则就只是模板。

\textbf{正确性。} 所有非零元素按原读取顺序写入前缀，因此稳定性保持；剩余位置数量等于原零的数量，补零后满足题意。

\textbf{边界实验。} 先用全零、无零、零在开头结尾、稳定顺序要求做反例检查。实验时把全零、无零、零在开头结尾、稳定顺序要求列成表，再打印关键下标和有效区间。若某个元素被写入后又被未来步骤破坏，说明区间不变量没有成立。

\textbf{复杂度变化。} 辅助数组 O(n) 空间；原地稳定写 O(n) 时间、O(1) 额外空间。

\textbf{迁移追问。} 如果题目改成“若目标是移动指定值，模型与移除元素相同”，先判断原状态是否仍然覆盖新答案集合，再决定是微调边界、改 value 语义，还是换算法。

\subsection{LeetCode 26 删除有序数组重复项}

\textbf{题目说明。} LeetCode 26 删除有序数组重复项 的题面入口要先还原成具体任务：有序数组让重复元素相邻，慢指针表示下一个唯一值写入位置。

\textbf{样例入口。} 拿 [1, 1, 2] 手算，write 先停在第一个 1 后面；读到第二个 1 时，它和已保留尾值相同，不写入；读到 2 时写到 write 位置，前缀变成 [1, 2]。

\textbf{暴力基线。} 暴力可以新建数组，把每组相同元素只写一次；或者每发现重复就移动后续元素覆盖它。

\textbf{暴力过程。} 以 [1, 1, 2] 为例，辅助数组写入 1，第二个 1 跳过，遇到 2 写入。原地做法只是把辅助数组的写入位置改成 nums[write]。

\textbf{重复工作。} 题面模型：有序数组让重复元素相邻，慢指针表示下一个唯一值写入位置。暴力版的慢点要落到本题：浪费在额外数组或反复搬移。数组已经有序，重复元素相邻，只需要和已保留前缀的最后一个值比较。后面的优化只消掉这类重复工作，不能改变答案集合。

\textbf{优化条件。} 本题的关键观察是：快指针扫描，只有发现新值才写入慢指针并推进。这句话必须能翻译成可维护状态；如果题目条件改变后观察不再成立，就不能继续套原来的写法。

\textbf{相邻状态。} 规律是慢指针左侧始终是不重复前缀，快指针只在遇到新值时推进慢指针；空数组和全重复数组只改变初始化，不改变不变量。把这个特例继续拆开看：先标出已经确认的前缀或后缀，再标出未知区；能被丢掉的候选，必须是以后再也不会改变已确认区域语义的候选。这样才算从例子推出数组不变量，而不是只记一次操作。

\textbf{状态复用。} write 表示下一个唯一值写入位置，read 从左到右扫描。nums[0..write-1] 始终是不重复前缀。手算时只改被新输入影响的那一小部分，而不是回头重算完整候选。

\textbf{指针和字段。} 状态字段要能说清：read 是当前读入位置，write 是有效前缀长度，nums[write-1] 是当前已保留的最后一个唯一值。手算时按这个协议跟踪：[1, 1, 2] 中 write 初始为 1；read=1 时 nums[1]==nums[0] 不写；read=2 时不同，写到 nums[1]，write 变 2。

\textbf{代码落点。} 关键代码是 if (nums[read] != nums[write-1]) nums[write++] = nums[read]; 。空数组直接返回 0。关键代码行必须能逐句翻译成“旧状态怎样变成新状态”，否则就只是模板。

\textbf{正确性。} 因为数组有序，同一个值只会连续出现；第一次出现被写入后，后续相同值都应跳过，遇到新值就扩展不重复前缀。

\textbf{边界实验。} 先用空数组、全重复、无重复、返回长度不等于物理容量做反例检查。实验时把空数组、全重复、无重复、返回长度不等于物理容量列成表，再打印关键下标和有效区间。若某个元素被写入后又被未来步骤破坏，说明区间不变量没有成立。

\textbf{复杂度变化。} 移动元素暴力最坏 O(\ensuremath{n^{2}}) 或辅助数组 O(n) 空间；快慢指针 O(n) 时间、O(1) 额外空间。

\textbf{迁移追问。} 如果题目改成“若允许每个元素最多保留两次，只需改变写入条件为和前两个比较”，先判断原状态是否仍然覆盖新答案集合，再决定是微调边界、改 value 语义，还是换算法。

\subsection{LeetCode 88 合并两个有序数组}

\textbf{题目说明。} LeetCode 88 合并两个有序数组 的题面入口要先还原成具体任务：第一个数组尾部有空位，适合从后往前写入最终最大值。

\textbf{样例入口。} 拿 nums1=[1, 3, 0]、nums2=[2] 手算，如果从前往后写，2 会覆盖 nums1 中尚未比较的 3；从尾部写，先比较 3 和 2，把 3 放到最后，再把 2 放到中间。

\textbf{暴力基线。} 暴力可以把 nums2 追加到 nums1 有效元素后再排序，或者新建数组从前往后归并。

\textbf{暴力过程。} 以 nums1=[1, 3, 0]、nums2=[2] 为例，从前往后写会把 nums1 中尚未读取的 3 覆盖掉。

\textbf{重复工作。} 题面模型：第一个数组尾部有空位，适合从后往前写入最终最大值。暴力版的慢点要落到本题：浪费在排序或额外数组。nums1 尾部本来有空位，应该从后往前写入当前最大值，避免覆盖未读元素。后面的优化只消掉这类重复工作，不能改变答案集合。

\textbf{优化条件。} 本题的关键观察是：逆向填充避免覆盖尚未读取的 nums1 元素。这句话必须能翻译成可维护状态；如果题目条件改变后观察不再成立，就不能继续套原来的写法。

\textbf{相邻状态。} 规律是利用尾部空位逆向填充，避免覆盖未读元素。把这个特例继续拆开看：先标出已经确认的前缀或后缀，再标出未知区；能被丢掉的候选，必须是以后再也不会改变已确认区域语义的候选。这样才算从例子推出数组不变量，而不是只记一次操作。

\textbf{状态复用。} i=m-1 指向 nums1 有效尾，j=n-1 指向 nums2 尾，write=m+n-1 指向最终写入位置。手算时只改被新输入影响的那一小部分，而不是回头重算完整候选。

\textbf{指针和字段。} 状态字段要能说清：i/j 是两个有序数组尚未合并部分的尾部，write 是结果数组尚未填充的最后位置。手算时按这个协议跟踪：[1, 3, 0] 和 [2] 中，比较 3 和 2，把 3 写到 write=2；再把 2 写到 write=1，1 已在原位。

\textbf{代码落点。} 关键代码是从尾部 while j>=0 合并；若 nums1 剩余不用搬，若 nums2 剩余必须继续写入。关键代码行必须能逐句翻译成“旧状态怎样变成新状态”，否则就只是模板。

\textbf{正确性。} 最终数组最大元素一定来自两个当前尾部之一，把它写到最后不会影响剩余元素的相对有序性。

\textbf{边界实验。} 先用第二个数组为空、第一个有效元素为空、重复值、尾部剩余做反例检查。实验时把第二个数组为空、第一个有效元素为空、重复值、尾部剩余列成表，再打印关键下标和有效区间。若某个元素被写入后又被未来步骤破坏，说明区间不变量没有成立。

\textbf{复杂度变化。} 排序 O((m+n)log(m+n))；逆向归并 O(m+n) 时间、O(1) 额外空间。

\textbf{迁移追问。} 如果题目改成“若没有额外空位，只能新开数组或使用更复杂原地归并”，先判断原状态是否仍然覆盖新答案集合，再决定是微调边界、改 value 语义，还是换算法。

\subsection{LeetCode 53 最大子数组和}

\textbf{题目说明。} LeetCode 53 最大子数组和给定整数数组 nums，要求找一个非空连续子数组，使元素和最大，并返回这个最大和。

\textbf{样例入口。} 样例 nums=[-2, 1, -3, 4, -1, 2, 1, -5, 4]，输出 6；连续段 [4, -1, 2, 1] 的和为 6。

\textbf{暴力基线。} 暴力枚举所有连续子数组，累计每段和并更新最大值；优化一点可以固定左端、右端扩展时维护当前和。

\textbf{暴力过程。} 以 [-2, 1, -3, 4, -1, 2, 1] 为例，暴力会重复计算很多以 4 开头或以 2 结尾的子段。真正关键是当前数要不要接上前一个后缀。

\textbf{重复工作。} 题面模型：给定整数数组，返回一个非空连续子数组能够取得的最大元素和。暴力版的慢点要落到本题：浪费在重复求每个后缀和。对于以 i 结尾的最优连续段，历史信息只需要一个值：以 i-1 结尾的最大后缀和。后面的优化只消掉这类重复工作，不能改变答案集合。

\textbf{优化条件。} 本题的关键观察是：用 [-2, 3] 和 [2, 3] 对比，推出当前位置只在单独开始与接上旧后缀之间取较大值。这句话必须能翻译成可维护状态；如果题目条件改变后观察不再成立，就不能继续套原来的写法。

\textbf{相邻状态。} 规律是当前位置只需要比较 nums[i] 与 bestEndingHere + nums[i]，全局答案再和 bestEndingHere 比较；全负数组要求答案初始化为第一个数，不能初始化为 0。把这个特例继续拆开看：先问暴力搜索里哪些不同路径到达同一个后续问题，再给这个后续问题命名为状态。状态必须包含未来决策需要的全部信息，转移只从更小且已经算清的状态来。

\textbf{状态复用。} bestEndingHere 表示以当前位置结尾的最大子数组和；answer 表示全局最大值。当前位置要么单独开始，要么接上旧后缀。手算时只改被新输入影响的那一小部分，而不是回头重算完整候选。

\textbf{指针和字段。} 状态字段要能说清：current=bestEndingHere，answer 是全局最优；nums[i] 与 current+nums[i] 比较决定是否丢弃旧后缀。手算时按这个协议跟踪：遇到 4 之前，旧后缀为负，接上反而变差，所以从 4 重新开始；之后 -1、2、1 继续接上得到 6。

\textbf{代码落点。} 关键代码是 current=max(nums[i], current+nums[i]); answer=max(answer, current)。answer 必须初始化为 nums[0]，否则全负数组会错。关键代码行必须能逐句翻译成“旧状态怎样变成新状态”，否则就只是模板。

\textbf{正确性。} 任意以 i 结尾的连续段只有两类：只含 nums[i]，或由某个以 i-1 结尾的连续段加上 nums[i]。取前一类最优即可覆盖全部候选。

\textbf{边界实验。} 先用全负数组、单元素、和溢出、答案不能初始化为零做反例检查。实验时覆盖全负数组、单元素、和溢出、答案不能初始化为零，先打印完整 DP 表，再比较空间压缩版本。若压缩版不同，优先检查遍历顺序和旧值保存。

\textbf{复杂度变化。} 暴力 O(\ensuremath{n^{2}}) 或 O(\ensuremath{n^{3}})，Kadane 算法 O(n) 时间、O(1) 空间。

\textbf{迁移追问。} 如果题目改成“若要求返回区间下标，同步记录重新开始位置和最优左右端”，先判断原状态是否仍然覆盖新答案集合，再决定是微调边界、改 value 语义，还是换算法。

\subsection{LeetCode 344 反转字符串}

\textbf{题目说明。} LeetCode 344 反转字符串 的题面入口要先还原成具体任务：反转只需要把对称位置交换，没必要新建完整副本。

\textbf{样例入口。} 拿 s=[h, e, l, l, o] 手算，交换首尾得到 o, e, l, l, h，再交换 e 和 l，最终 o, l, l, e, h。

\textbf{暴力基线。} 暴力可以新建一个反向数组，把原数组从后往前复制过去，再写回。

\textbf{暴力过程。} 以 [h, e, l, l, o] 为例，辅助数组得到 [o, l, l, e, h]。原地做法只需要交换首尾、次首尾。

\textbf{重复工作。} 题面模型：反转只需要把对称位置交换，没必要新建完整副本。暴力版的慢点要落到本题：浪费在完整副本。反转的每个位置只和它的对称位置交换一次。后面的优化只消掉这类重复工作，不能改变答案集合。

\textbf{优化条件。} 本题的关键观察是：左右指针从两端向中间移动，每轮交换一对字符。这句话必须能翻译成可维护状态；如果题目条件改变后观察不再成立，就不能继续套原来的写法。

\textbf{相邻状态。} 规律是左右指针每轮交换对称字符并向中间收缩，原地反转不需要额外数组。把这个特例继续拆开看：先标出已经确认的前缀或后缀，再标出未知区；能被丢掉的候选，必须是以后再也不会改变已确认区域语义的候选。这样才算从例子推出数组不变量，而不是只记一次操作。

\textbf{状态复用。} left 从 0 开始，right 从 n-1 开始；left<right 时交换 s[left] 和 s[right]，然后向中间收缩。手算时只改被新输入影响的那一小部分，而不是回头重算完整候选。

\textbf{指针和字段。} 状态字段要能说清：left/right 是尚未交换的对称位置，二者之间是还未完全处理的区域。手算时按这个协议跟踪：先交换 h 和 o，得到 [o, e, l, l, h]；再交换 e 和 l，得到 [o, l, l, e, h]。

\textbf{代码落点。} 关键代码是 while (left < right) swap(s[left++], s[right--])。关键代码行必须能逐句翻译成“旧状态怎样变成新状态”，否则就只是模板。

\textbf{正确性。} 每次交换把一对对称位置放到最终位置，已处理两端不会再被修改，未知区逐步缩小。

\textbf{边界实验。} 先用空数组、单字符、偶数长度、奇数长度、原地修改做反例检查。实验时把空数组、单字符、偶数长度、奇数长度、原地修改列成表，再打印关键下标和有效区间。若某个元素被写入后又被未来步骤破坏，说明区间不变量没有成立。

\textbf{复杂度变化。} 辅助数组 O(n) 空间；双指针 O(n) 时间、O(1) 额外空间。

\textbf{迁移追问。} 如果题目改成“若反转单词顺序，要先定义分隔符和单词边界”，先判断原状态是否仍然覆盖新答案集合，再决定是微调边界、改 value 语义，还是换算法。

\subsection{LeetCode 125 验证回文串}

\textbf{题目说明。} LeetCode 125 验证回文串 的题面入口要先还原成具体任务：忽略非字母数字后，回文只需要两端字符向中间比较。

\textbf{样例入口。} 拿 A man, a plan, a canal: Panama 手算，忽略空格和标点后得到 amanaplanacanalpanama，左右两端逐个比较都相同。

\textbf{暴力基线。} 暴力可以先构造一个只含字母数字且统一小写的新字符串，再反转或双指针比较。

\textbf{暴力过程。} 以 A man, a plan, a canal: Panama 为例，过滤后是 amanaplanacanalpanama，再判断它是否回文。

\textbf{重复工作。} 题面模型：忽略非字母数字后，回文只需要两端字符向中间比较。暴力版的慢点要落到本题：浪费在新建过滤字符串。原字符串两端向中间走时，可以边跳过无关字符边比较。后面的优化只消掉这类重复工作，不能改变答案集合。

\textbf{优化条件。} 本题的关键观察是：左右指针分别跳过非法字符，再比较小写化后的字符，相等则同时收缩。这句话必须能翻译成可维护状态；如果题目条件改变后观察不再成立，就不能继续套原来的写法。

\textbf{相邻状态。} 规律是左右指针跳过非字母数字字符，比较时统一大小写，指针只向中间移动。把这个特例继续拆开看：右端进入只带来一个新元素，左端离开只删掉一个旧元素；只有当旧左端被移走后不可能再产生更优答案时，窗口才可以单向推进。若这个单调淘汰条件不存在，就要换成前缀和、队列或其他状态。

\textbf{状态复用。} left/right 从两端开始，跳过非字母数字字符，比较 tolower 后的字符；相等则同时收缩。手算时只改被新输入影响的那一小部分，而不是回头重算完整候选。

\textbf{指针和字段。} 状态字段要能说清：left/right 是尚未确认的两端字符位置，合法字符判断决定是否参与比较。手算时按这个协议跟踪：左端 A 和右端 a 统一小写后相等；中间遇到空格、逗号、冒号都跳过。

\textbf{代码落点。} 关键是调用 isalnum 和 tolower 时用 unsigned char 转换，避免 char 为负时行为不稳；比较失败立即返回 false。关键代码行必须能逐句翻译成“旧状态怎样变成新状态”，否则就只是模板。

\textbf{正确性。} 回文只要求过滤后的第 k 个字符等于倒数第 k 个字符；双指针跳过无关字符后比较的正是这两者。

\textbf{边界实验。} 先用空串、只有标点、大小写、数字字符、Unicode 范围假设做反例检查。实验时重点跑空串、只有标点、大小写、数字字符、Unicode 范围假设，并打印左右端、窗口计数和答案。若左端发生回退，或者窗口失去合法性后仍更新答案，就能很快定位错误。

\textbf{复杂度变化。} 构造法 O(n) 时间和 O(n) 空间；原地双指针 O(n) 时间、O(1) 空间。

\textbf{迁移追问。} 如果题目改成“若要求最多删除一个字符，可以在第一次失配时分支检查两个子区间”，先判断原状态是否仍然覆盖新答案集合，再决定是微调边界、改 value 语义，还是换算法。

\subsection{LeetCode 11 盛最多水的容器}

\textbf{题目说明。} LeetCode 11 盛最多水的容器给定数组 height，height[i] 表示第 i 条竖线高度。选择两条线与 x 轴形成容器，返回能容纳的最大水量。

\textbf{样例入口。} 样例 height=[1, 8, 6, 2, 5, 4, 8, 3, 7]，输出 49；选择下标 1 和 8，宽度 7，短板高度 7，面积 49。

\textbf{暴力基线。} 暴力枚举所有 left<right 的下标对，面积等于 (right-left)*min(height[left], height[right])，取最大值。

\textbf{暴力过程。} 以 [1, 8, 6, 2, 5, 4, 8, 3, 7] 为例，暴力会试 (0, 1)、(0, 2) 一直到 (0, 8)，再试 (1, 2)、(1, 3)。大量候选只是宽度和短板变化。

\textbf{重复工作。} 题面模型：给定每个下标的高度，选择两条竖线和横轴组成容器，返回能盛水的最大面积。暴力版的慢点要落到本题：浪费在逐个试所有右端。若当前短板是 left，保留 left 再把 right 左移只会让宽度变小，短板高度不可能超过 height[left]，这些候选都被当前状态支配。后面的优化只消掉这类重复工作，不能改变答案集合。

\textbf{优化条件。} 本题的关键观察是：每次移动短板，因为移动长板只会缩小宽度且不会提高短板。这句话必须能翻译成可维护状态；如果题目条件改变后观察不再成立，就不能继续套原来的写法。

\textbf{相邻状态。} 规律是每轮淘汰短板那一侧，因为所有保留它且宽度更小的候选都被当前面积上界支配。把这个特例继续拆开看：右端进入只带来一个新元素，左端离开只删掉一个旧元素；只有当旧左端被移走后不可能再产生更优答案时，窗口才可以单向推进。若这个单调淘汰条件不存在，就要换成前缀和、队列或其他状态。

\textbf{状态复用。} 双指针从两端开始，维护 left、right 和 best。每轮计算当前面积，然后移动高度较小的一侧。手算时只改被新输入影响的那一小部分，而不是回头重算完整候选。

\textbf{指针和字段。} 状态字段要能说清：left/right 是当前最宽候选的两端，shorter 是限制高度的一侧，best 是已见过最大面积。手算时按这个协议跟踪：样例中 left=0 高 1、right=8 高 7，面积 8。短板是 left，移动 left 到 1；此时高 8 和 7，面积 49，更新答案。

\textbf{代码落点。} 关键代码是 if (height[left] < height[right]) ++left; else --right; 。移动长板无法提高短板上界，所以不移动长板。关键代码行必须能逐句翻译成“旧状态怎样变成新状态”，否则就只是模板。

\textbf{正确性。} 每次移动短板时，所有包含这个短板且宽度更小的候选面积都不可能超过当前面积上界，因此可以一次排除一批候选。

\textbf{边界实验。} 先用两端高度相等、数组长度为二、面积乘法可能溢出做反例检查。实验时重点跑两端高度相等、数组长度为二、面积乘法可能溢出，并打印左右端、窗口计数和答案。若左端发生回退，或者窗口失去合法性后仍更新答案，就能很快定位错误。

\textbf{复杂度变化。} 暴力 O(\ensuremath{n^{2}})，双指针每轮移动一端，总共 O(n) 轮，时间 O(n)，空间 O(1)。

\textbf{迁移追问。} 如果题目改成“接雨水计算每个位置贡献，不能把本题双指针证明直接搬过去”，先判断原状态是否仍然覆盖新答案集合，再决定是微调边界、改 value 语义，还是换算法。

\subsection{LeetCode 15 三数之和}

\textbf{题目说明。} LeetCode 15 三数之和给定整数数组 nums，要求找出所有不重复三元组 [a, b, c]，使 a+b+c=0。

\textbf{样例入口。} 样例 nums=[-1, 0, 1, 2, -1, -4]，输出 [[-1, -1, 2], [-1, 0, 1]]；结果里不能出现重复三元组。

\textbf{暴力基线。} 暴力三层循环枚举 i、j、k，检查 nums[i]+nums[j]+nums[k] 是否为 0，再用集合去掉重复三元组。

\textbf{暴力过程。} 以 [-1, 0, 1, 2, -1, -4] 为例，暴力会枚举 (-1, 0, 1)、(-1, 0, 2)、(-1, 0, -1) 等所有三元组；两个 -1 会制造重复答案。

\textbf{重复工作。} 题面模型：给定整数数组，返回所有不重复三元组，使三个数之和等于零。暴力版的慢点要落到本题：浪费在第三层枚举和事后去重。排序后固定一个数，剩余两个数变成有序两数之和，可以根据当前和一次排除一侧。后面的优化只消掉这类重复工作，不能改变答案集合。

\textbf{优化条件。} 本题的关键观察是：排序后固定第一个数，剩余两数转成有序双指针；固定值和左右值都要跳过重复。这句话必须能翻译成可维护状态；如果题目条件改变后观察不再成立，就不能继续套原来的写法。

\textbf{相邻状态。} 规律是排序把三数问题拆成枚举固定点加双指针，并在每层处理去重。把这个特例继续拆开看：右端进入只带来一个新元素，左端离开只删掉一个旧元素；只有当旧左端被移走后不可能再产生更优答案时，窗口才可以单向推进。若这个单调淘汰条件不存在，就要换成前缀和、队列或其他状态。

\textbf{状态复用。} 先排序。外层 i 固定第一个数，内层 left=i+1、right=n-1 搜索 target=-nums[i]；找到答案后跳过相同 left/right 值。手算时只改被新输入影响的那一小部分，而不是回头重算完整候选。

\textbf{指针和字段。} 状态字段要能说清：i 是本层固定点，left/right 是右侧有序区间的候选两端，sum 是三数和，res 保存已去重结果。手算时按这个协议跟踪：排序后 [-4, -1, -1, 0, 1, 2]。i=1 固定 -1，left 指向 -1，right 指向 2，和为 0，记录 [-1, -1, 2]，然后跳过重复。

\textbf{代码落点。} 关键代码是 sum<0 时 ++left，sum>0 时 --right；i>0 且 nums[i]==nums[i-1] 要 continue，命中后 left/right 也要跳过重复值。关键代码行必须能逐句翻译成“旧状态怎样变成新状态”，否则就只是模板。

\textbf{正确性。} 排序后，当前和太小，固定 i 和 right 时更小的 left 都不可能达标；当前和太大，固定 i 和 left 时更大的 right 都不可能达标。去重发生在同一层，避免漏掉不同层组合。

\textbf{边界实验。} 先用全零、重复元素很多、没有答案、结果顺序不要求但组合不能重复做反例检查。实验时重点跑全零、重复元素很多、没有答案、结果顺序不要求但组合不能重复，并打印左右端、窗口计数和答案。若左端发生回退，或者窗口失去合法性后仍更新答案，就能很快定位错误。

\textbf{复杂度变化。} 暴力 O(\ensuremath{n^{3}}) 加去重开销；排序 O(n log n)，外层 n 次、内层双指针线性，总时间 O(\ensuremath{n^{2}})，输出空间另计。

\textbf{迁移追问。} 如果题目改成“四数之和再外层固定一个数，并用 long long 防止目标和溢出”，先判断原状态是否仍然覆盖新答案集合，再决定是微调边界、改 value 语义，还是换算法。

\subsection{LeetCode 209 长度最小的子数组}

\textbf{题目说明。} LeetCode 209 长度最小的子数组给定正整数数组 nums 和目标 target，返回和至少为 target 的最短连续子数组长度；不存在则返回 0。

\textbf{样例入口。} 样例 target=7，nums=[2, 3, 1, 2, 4, 3]，输出 2；最短合法连续段是 [4, 3]。

\textbf{暴力基线。} 暴力枚举每个起点 left，向右累计 sum，第一次 sum>=target 时更新长度，再换下一个起点。

\textbf{暴力过程。} 以 target=7、nums=[2, 3, 1, 2, 4, 3] 为例，从 0 开始要扩到 [2, 3, 1, 2]；从 1 开始又重新累加 [3, 1, 2, 4]。

\textbf{重复工作。} 题面模型：给定正整数数组和目标值，返回和至少为目标值的最短连续子数组长度。暴力版的慢点要落到本题：浪费在相邻起点重复累加。因为 nums 全为正数，右端增加只会让 sum 变大，左端移走只会让 sum 变小。后面的优化只消掉这类重复工作，不能改变答案集合。

\textbf{优化条件。} 本题的关键观察是：正数数组让窗口和随右端增加、随左端减少；右端扩展到满足目标后，不断左移缩短并更新答案。这句话必须能翻译成可维护状态；如果题目条件改变后观察不再成立，就不能继续套原来的写法。

\textbf{相邻状态。} 规律是全正数保证右扩和变大、左缩和变小，所以每个左端只需被移动一次。把这个特例继续拆开看：右端进入只带来一个新元素，左端离开只删掉一个旧元素；只有当旧左端被移走后不可能再产生更优答案时，窗口才可以单向推进。若这个单调淘汰条件不存在，就要换成前缀和、队列或其他状态。

\textbf{状态复用。} left、right 维护当前窗口和 sum。right 每次加入一个数；当 sum>=target 时，循环收缩 left 并更新最短长度。手算时只改被新输入影响的那一小部分，而不是回头重算完整候选。

\textbf{指针和字段。} 状态字段要能说清：sum 是 nums[left..right] 的和，answer 是已见过最短合法长度；left 只在窗口合法时右移。手算时按这个协议跟踪：样例中 right 到下标 3 时 sum=8，记录长度 4，移走 2 后 sum=6 停止；后来窗口 [4, 3] sum=7，答案更新为 2。

\textbf{代码落点。} 关键代码是 sum += nums[right]; while (sum >= target) \{ ans=min(ans, right-left+1); sum-=nums[left++]; \}。关键代码行必须能逐句翻译成“旧状态怎样变成新状态”，否则就只是模板。

\textbf{正确性。} 正数保证收缩 left 后 sum 单调下降；对每个 right，while 会找到以它结尾的所有可缩短合法窗口，不会漏最短值。

\textbf{边界实验。} 先用不存在答案、单元素达到目标、全正是必要条件做反例检查。实验时重点跑不存在答案、单元素达到目标、全正是必要条件，并打印左右端、窗口计数和答案。若左端发生回退，或者窗口失去合法性后仍更新答案，就能很快定位错误。

\textbf{复杂度变化。} 暴力 O(\ensuremath{n^{2}})，窗口中每个元素至多进入一次、离开一次，时间 O(n)，空间 O(1)。

\textbf{迁移追问。} 如果题目改成“若含负数，需要前缀和加单调队列”，先判断原状态是否仍然覆盖新答案集合，再决定是微调边界、改 value 语义，还是换算法。

\subsection{LeetCode 3 无重复字符的最长子串}

\textbf{题目说明。} LeetCode 3 无重复字符的最长子串 的题面入口要先还原成具体任务：窗口保存当前无重复子串，右端每次加入一个字符。

\textbf{样例入口。} 拿字符串 abca 手算。窗口 abc 合法，右端加入第二个 a 后，真正冲突的是上一次 a 的位置；左端从 0 跳到 1 后，bca 重新合法。再试 abba：处理第二个 b 时左端跳到 2，后面遇到 a 的上次位置在窗口左侧之外，不能让 left 回退。

\textbf{暴力基线。} 暴力固定每个起点 left，向右逐个加入字符，用集合检查是否重复；一旦重复就停止当前起点，换下一个起点重新建集合。

\textbf{暴力过程。} 以 s=abcabcbb 为例，暴力会检查 a、ab、abc，到 abca 时发现 a 重复；再从 b 重新检查 b、bc、bca。相邻起点的窗口 b、c 其实已经在上一轮出现过。

\textbf{重复工作。} 题面模型：窗口保存当前无重复子串，右端每次加入一个字符。暴力版的慢点要落到本题：浪费在每换一个 left 就重新统计一段已经看过的字符。真正变化只是左端离开若干旧字符，右端继续加入一个新字符。后面的优化只消掉这类重复工作，不能改变答案集合。

\textbf{优化条件。} 本题的关键观察是：用上次出现位置直接跳过无效左端，左边界只能向右不能回退。这句话必须能翻译成可维护状态；如果题目条件改变后观察不再成立，就不能继续套原来的写法。

\textbf{相邻状态。} 规律是左边界只向右，状态保存每个字符最近出现位置。把这个特例继续拆开看：右端进入只带来一个新元素，左端离开只删掉一个旧元素；只有当旧左端被移走后不可能再产生更优答案时，窗口才可以单向推进。若这个单调淘汰条件不存在，就要换成前缀和、队列或其他状态。

\textbf{状态复用。} 保存窗口左端 left、右端 right、每个字符最近出现位置 last、当前最长答案 best。遇到重复字符时，left 直接跳到 last[ch]+1，但只能向右跳。手算时只改被新输入影响的那一小部分，而不是回头重算完整候选。

\textbf{指针和字段。} 状态字段要能说清：right 是当前加入字符位置，left 是当前无重复窗口起点，last[ch] 是字符 ch 上一次出现下标，best 是已经见过的最大合法窗口长度。手算时按这个协议跟踪：手算 abba：right=0 加 a，left=0；right=1 加 b；right=2 再遇 b，left 跳到 2；right=3 遇 a，但上次 a 在 0，已经在窗口外，left 不能回退。

\textbf{代码落点。} 关键代码是 left=max(left, last[ch]+1)，然后 last[ch]=right，最后用 right-left+1 更新答案。顺序不能反，否则会把当前字符当成自己的历史位置。关键代码行必须能逐句翻译成“旧状态怎样变成新状态”，否则就只是模板。

\textbf{正确性。} left 跳过的是所有仍包含重复 ch 的起点，这些起点不可能形成以 right 结尾的合法窗口；left 不回退保证每个字符至多进出窗口一次。

\textbf{边界实验。} 先用空串、全相同字符、重复字符出现在窗口左侧之外、扩展字符集做反例检查。实验时重点跑空串、全相同字符、重复字符出现在窗口左侧之外、扩展字符集，并打印左右端、窗口计数和答案。若左端发生回退，或者窗口失去合法性后仍更新答案，就能很快定位错误。

\textbf{复杂度变化。} 暴力最坏 O(\ensuremath{n^{2}})，优化后每个 right 扫一次，哈希或数组访问均摊 O(1)，总时间 O(n)，空间取决于字符集。

\textbf{迁移追问。} 如果题目改成“若要求恰好包含 k 种字符，窗口条件从无重复变成种类计数”，先判断原状态是否仍然覆盖新答案集合，再决定是微调边界、改 value 语义，还是换算法。

\subsection{LeetCode 438 找到字符串中所有字母异位词}

\textbf{题目说明。} LeetCode 438 找到字符串中所有字母异位词 的题面入口要先还原成具体任务：固定长度窗口内字符频次等于目标频次即为答案。

\textbf{样例入口。} 拿 s=cbaebabacd、p=abc 手算，长度为 3 的窗口 cba 频次和 abc 相同，起点 0 是答案；窗口右移时只删除 c、加入 e。

\textbf{暴力基线。} 先统计 p 的频次 target，再枚举 s 中每个长度 k=p.size() 的窗口，为每个窗口重新统计 window 后比较。

\textbf{暴力过程。} 以 s=cbaebabacd、p=abc 为例，暴力窗口依次是 cba、bae、aeb、eba、bab、aba、bac、acd。每个窗口都重新数 k 个字符。

\textbf{重复工作。} 题面模型：固定长度窗口内字符频次等于目标频次即为答案。暴力版的慢点要落到本题：相邻窗口只差一个出去字符和一个进来字符。从 cba 到 bae，只有 c 出去、e 进来，b 和 a 不该重数。后面的优化只消掉这类重复工作，不能改变答案集合。

\textbf{优化条件。} 本题的关键观察是：右进左出维护计数，窗口长度达到目标长度后检查差异。这句话必须能翻译成可维护状态；如果题目条件改变后观察不再成立，就不能继续套原来的写法。

\textbf{相邻状态。} 规律是固定长度窗口维护字符频次差，重叠答案不能漏。把这个特例继续拆开看：右端进入只带来一个新元素，左端离开只删掉一个旧元素；只有当旧左端被移走后不可能再产生更优答案时，窗口才可以单向推进。若这个单调淘汰条件不存在，就要换成前缀和、队列或其他状态。

\textbf{状态复用。} k 来自 p.length()。先统计第一个窗口，之后每次窗口右移一格：window[s[i-1]]--，window[s[i+k-1]]++。手算时只改被新输入影响的那一小部分，而不是回头重算完整候选。

\textbf{指针和字段。} 状态字段要能说清：i 是当前窗口起点，out=s[i-1] 是离开字符，in=s[i+k-1] 是进入字符，target/window 是 26 长度频次数组。手算时按这个协议跟踪：第一个窗口 cba 与 abc 频次相同，记录 0；滑到 bae 后 c 计数减一、e 加一；滑到 bac 时频次再次相同，记录 6。

\textbf{代码落点。} 关键代码是固定窗口右进左出。若用 unordered\_map，计数减到 0 的 key 要 erase；若题目限定小写字母，用 array/vector 长度 26 更简单。关键代码行必须能逐句翻译成“旧状态怎样变成新状态”，否则就只是模板。

\textbf{正确性。} 每个被检查窗口长度都等于 k，window 频次始终精确表示当前窗口；频次等于 target 当且仅当该窗口是 p 的异位词。

\textbf{边界实验。} 先用目标比源串长、重复字符、字符集固定、答案重叠做反例检查。实验时重点跑目标比源串长、重复字符、字符集固定、答案重叠，并打印左右端、窗口计数和答案。若左端发生回退，或者窗口失去合法性后仍更新答案，就能很快定位错误。

\textbf{复杂度变化。} 暴力 O(nk)，固定窗口维护 O(n*26)，因字符集固定通常记作 O(n)，空间 O(26)。

\textbf{迁移追问。} 如果题目改成“维护差异计数可以避免每次比较整个数组”，先判断原状态是否仍然覆盖新答案集合，再决定是微调边界、改 value 语义，还是换算法。

\subsection{LeetCode 76 最小覆盖子串}

\textbf{题目说明。} LeetCode 76 最小覆盖子串给定字符串 s 和 t，要求在 s 中找最短连续子串，使它包含 t 中每个字符及其次数；不存在则返回空串。

\textbf{样例入口。} 样例 s=ADOBECODEBANC，t=ABC，输出 BANC；它是覆盖 A、B、C 的最短连续子串。

\textbf{暴力基线。} 暴力枚举 s 的每个子串，重新统计字符频次，检查是否覆盖 t 的每个字符及其次数。

\textbf{暴力过程。} 以 s=ADOBECODEBANC、t=ABC 为例，暴力会检查 A、AD、ADO...直到 ADOBEC 才第一次覆盖；下一轮从 D 开始又重新统计 DOBEC。

\textbf{重复工作。} 题面模型：给定字符串 s 和 t，在 s 中找一个最短连续子串，使它覆盖 t 的每个字符及其次数。暴力版的慢点要落到本题：浪费在相邻子串重复统计字符。右端扩展只新增一个字符，左端收缩只移走一个字符。后面的优化只消掉这类重复工作，不能改变答案集合。

\textbf{优化条件。} 本题的关键观察是：窗口内字符计数必须覆盖目标计数，满足后尽量收缩左端；用 formed 记录满足需求的字符种类，避免每次全量比较。这句话必须能翻译成可维护状态；如果题目条件改变后观察不再成立，就不能继续套原来的写法。

\textbf{相邻状态。} 规律是右端负责获得可行，左端负责在可行时尽量缩短；计数表要区分已有字符和需求字符。把这个特例继续拆开看：右端进入只带来一个新元素，左端离开只删掉一个旧元素；只有当旧左端被移走后不可能再产生更优答案时，窗口才可以单向推进。若这个单调淘汰条件不存在，就要换成前缀和、队列或其他状态。

\textbf{状态复用。} need 保存目标频次，window 保存当前窗口频次，formed 或 required\_count 记录当前满足需求的字符种类；满足后尽量收缩 left。手算时只改被新输入影响的那一小部分，而不是回头重算完整候选。

\textbf{指针和字段。} 状态字段要能说清：right 负责扩展获得覆盖，left 负责在覆盖时缩短窗口，best\_left/best\_len 记录最短答案。手算时按这个协议跟踪：ADOBEC 第一次覆盖 ABC，记录长度 6；左端移走 A 后不覆盖，停止收缩。继续扩展到 BANC 时再次覆盖，收缩得到长度 4。

\textbf{代码落点。} 关键顺序是加入 s[right] 后更新 formed；while formed==required 时先更新答案，再移走 s[left] 并可能降低 formed。关键代码行必须能逐句翻译成“旧状态怎样变成新状态”，否则就只是模板。

\textbf{正确性。} 每个 right 下，while 循环会把 left 推到刚好不能再缩的位置，因此不会错过以该 right 结尾的最短覆盖窗口。

\textbf{边界实验。} 先用目标有重复字符、源串缺字符、最小窗口在开头或结尾、大小写做反例检查。实验时重点跑目标有重复字符、源串缺字符、最小窗口在开头或结尾、大小写，并打印左右端、窗口计数和答案。若左端发生回退，或者窗口失去合法性后仍更新答案，就能很快定位错误。

\textbf{复杂度变化。} 暴力至少 O(\ensuremath{n^{2}}*字符集)，滑动窗口每个字符进出一次，时间 O(n+字符集)，空间 O(字符集)。

\textbf{迁移追问。} 如果题目改成“若字符集固定，数组计数更快；若是 Unicode，哈希表更通用”，先判断原状态是否仍然覆盖新答案集合，再决定是微调边界、改 value 语义，还是换算法。

\subsection{LeetCode 739 每日温度}

\textbf{题目说明。} LeetCode 739 每日温度 的题面入口要先还原成具体任务：每一天等待右侧第一个更高温度。

\textbf{样例入口。} 拿 [73, 74, 75, 71, 69, 72, 76, 73] 手算，73 遇到 74 时等待 1 天；71 要等到 72 才解决。

\textbf{暴力基线。} 慢版本让每个位置向左或向右线性寻找边界，或者让每个结构反复等待匹配。它能算出答案，但很多尚未完成的候选会被未来同一个事件一起解决。放到本题就是：先按“每一天等待右侧第一个更高温度”完整试慢版本；样例里已经能看到“规律是单调递减栈保存尚未等到更高温的下标，遇到更高温时连续弹栈结算天数”，所以优化前必须先能说清慢版本枚举了哪些候选。

\textbf{暴力过程。} 拿 [73, 74, 75, 71, 69, 72, 76, 73] 手算，73 遇到 74 时等待 1 天；71 要等到 72 才解决。

\textbf{重复工作。} 题面模型：每一天等待右侧第一个更高温度。暴力版的慢点要落到本题：慢版本会围绕“每一天等待右侧第一个更高温度”反复寻找最近边界或最近未完成对象。样例规律是“规律是单调递减栈保存尚未等到更高温的下标，遇到更高温时连续弹栈结算天数”，说明旧候选一旦遇到当前输入就能被批量结算；继续为每个位置单独向左或向右扫描就是浪费。后面的优化只消掉这类重复工作，不能改变答案集合。

\textbf{优化条件。} 本题的关键观察是：单调递减栈保存尚未找到答案的下标。这句话必须能翻译成可维护状态；如果题目条件改变后观察不再成立，就不能继续套原来的写法。

\textbf{相邻状态。} 规律是单调递减栈保存尚未等到更高温的下标，遇到更高温时连续弹栈结算天数。把这个特例继续拆开看：栈里的对象都是答案尚未确定的候选；一旦当前元素触发弹栈，被弹出的候选必须已经同时拿到了左边界和右边界，未来元素不再有资格改变它的答案。

\textbf{状态复用。} 把反复寻找最近边界改成维护未完成候选；新元素只和栈顶比较，连续弹出一批已经被当前元素解决的对象。本题落点是：规律是单调递减栈保存尚未等到更高温的下标，遇到更高温时连续弹栈结算天数。手算时只改被新输入影响的那一小部分，而不是回头重算完整候选。

\textbf{指针和字段。} 状态字段要能说清：栈内对象的业务含义、当前输入、弹出时结算的对象、左边界和右边界；栈中存下标还是值要明确。手算时按这个协议跟踪：拿 [73, 74, 75, 71, 69, 72, 76, 73] 手算，73 遇到 74 时等待 1 天；71 要等到 72 才解决。然后按协议复查：手算时记录栈内元素的业务含义，不只记录数值。入栈表示答案未定，弹栈表示当前事件已经让它的答案定下来。

\textbf{代码落点。} 弹栈前确认栈非空，弹出后再读取新的左边界；相等元素和哨兵高度要有统一策略。关键代码行必须能逐句翻译成“旧状态怎样变成新状态”，否则就只是模板。

\textbf{正确性。} 证明从弹出时机开始：被弹出的对象在当前时刻答案已经确定，未来元素无法再改变它。

\textbf{边界实验。} 先用没有更高温、相等温度、递增、递减做反例检查。实验时准备没有更高温、相等温度、递增、递减，打印每次入栈、弹栈和结算对象。重点观察相等元素、空栈、哨兵和最后一次结算。

\textbf{复杂度变化。} 暴力为每个位置寻找边界常见为平方级；单调栈让每个元素最多入栈一次、出栈一次，因此主体扫描为线性时间。

\textbf{迁移追问。} 如果题目改成“下一个更大元素与它同型，只是输出值或距离不同”，先判断原状态是否仍然覆盖新答案集合，再决定是微调边界、改 value 语义，还是换算法。

\subsection{LeetCode 84 柱状图中最大的矩形}

\textbf{题目说明。} LeetCode 84 柱状图中最大的矩形 的题面入口要先还原成具体任务：每根柱子作为最矮高度时，需要知道左右第一个更矮位置。

\textbf{样例入口。} 拿高度 [2, 1, 2] 手算。第一个 2 遇到更矮的 1 时，右边界已经确定在当前位置，左边界是弹出后新栈顶的后一位，所以面积可结算。若一直递增，最后需要哨兵触发结算。

\textbf{暴力基线。} 暴力枚举每根柱子作为矩形最矮高度，再向左右扩展直到遇到更矮柱子，计算最大面积。

\textbf{暴力过程。} 以 heights=[2, 1, 5, 6, 2, 3] 为例，高度 5 的柱子向右可覆盖 6，向左遇到 1 停止，面积可为 5*2=10。

\textbf{重复工作。} 题面模型：每根柱子作为最矮高度时，需要知道左右第一个更矮位置。暴力版的慢点要落到本题：浪费在每根柱子都重新找左右第一个更矮位置。单调递增栈能在遇到更矮柱时一次结算被弹出柱子的右边界。后面的优化只消掉这类重复工作，不能改变答案集合。

\textbf{优化条件。} 本题的关键观察是：单调递增栈在遇到更矮柱时结算被弹出柱子的最大宽度。这句话必须能翻译成可维护状态；如果题目条件改变后观察不再成立，就不能继续套原来的写法。

\textbf{相邻状态。} 规律是栈保存递增高度下标，弹出时才知道该柱子的最大扩展宽度。把这个特例继续拆开看：栈里的对象都是答案尚未确定的候选；一旦当前元素触发弹栈，被弹出的候选必须已经同时拿到了左边界和右边界，未来元素不再有资格改变它的答案。

\textbf{状态复用。} 栈保存下标，且高度单调递增。扫描到更矮高度时，弹出 top，当前下标是右边界，弹出后新的栈顶是左边界。手算时只改被新输入影响的那一小部分，而不是回头重算完整候选。

\textbf{指针和字段。} 状态字段要能说清：top 是被结算柱子，height[top] 是矩形高度，left\_smaller=stack.top，right\_smaller=i，宽度为 i-left\_smaller-1。手算时按这个协议跟踪：扫描到高度 2 时，6 和 5 依次被弹出；弹出 5 时左边界是高度 1 的下标，右边界是当前 2 的下标，面积 10。

\textbf{代码落点。} 关键是两端加哨兵 0，或循环结束后手动清栈。相等高度要有统一策略，常见写法用 >= 弹出保证边界一致。关键代码行必须能逐句翻译成“旧状态怎样变成新状态”，否则就只是模板。

\textbf{正确性。} 柱子被弹出时，左右第一个更矮位置都已经确定，未来更右的柱子不能再扩大以它为最低高度的矩形。

\textbf{边界实验。} 先用递增、递减、相等高度、哨兵零、宽度计算做反例检查。实验时准备递增、递减、相等高度、哨兵零、宽度计算，打印每次入栈、弹栈和结算对象。重点观察相等元素、空栈、哨兵和最后一次结算。

\textbf{复杂度变化。} 暴力 O(\ensuremath{n^{2}})，单调栈每个下标入栈出栈一次，O(n) 时间、O(n) 空间。

\textbf{迁移追问。} 如果题目改成“最大矩形可以作为二维矩阵最大矩形的行压缩子问题”，先判断原状态是否仍然覆盖新答案集合，再决定是微调边界、改 value 语义，还是换算法。

\subsection{LeetCode 239 滑动窗口最大值}

\textbf{题目说明。} LeetCode 239 滑动窗口最大值 的题面入口要先还原成具体任务：每个窗口最大值来自仍未过期且没有被新元素支配的候选。

\textbf{样例入口。} 拿 [1, 3, -1]、窗口 3 手算，队尾的 1 在 3 进入后永远不可能成为最大值，因为 3 更新且更靠右。于是 1 可以弹出。

\textbf{暴力基线。} 暴力枚举每个长度为 k 的窗口，扫描窗口内所有元素求最大值。

\textbf{暴力过程。} 以 [1, 3, -1, -3, 5]、k=3 为例，第一个窗口 [1, 3, -1] 最大值为 3；下一个窗口 [3, -1, -3] 仍要重新扫描会重复看 3 和 -1。

\textbf{重复工作。} 题面模型：每个窗口最大值来自仍未过期且没有被新元素支配的候选。暴力版的慢点要落到本题：浪费在相邻窗口重复求最大值。若新来的数更大且位置更靠右，队尾所有不大于它的旧数未来都不可能再当最大值。后面的优化只消掉这类重复工作，不能改变答案集合。

\textbf{优化条件。} 本题的关键观察是：单调队列保存下标，队尾小于等于新值的候选被弹出。这句话必须能翻译成可维护状态；如果题目条件改变后观察不再成立，就不能继续套原来的写法。

\textbf{相邻状态。} 规律是单调队列保存可能成为当前或未来窗口最大值的下标，过期下标从队首清理。把这个特例继续拆开看：栈里的对象都是答案尚未确定的候选；一旦当前元素触发弹栈，被弹出的候选必须已经同时拿到了左边界和右边界，未来元素不再有资格改变它的答案。

\textbf{状态复用。} 单调递减队列保存下标。入队前从队尾弹出 nums[back]<=nums[i] 的下标；队首若过期 i-k 则弹出；窗口形成后队首就是最大值。手算时只改被新输入影响的那一小部分，而不是回头重算完整候选。

\textbf{指针和字段。} 状态字段要能说清：deque.front 是当前最大值下标，i 是右端，i-k+1 是当前窗口左端。手算时按这个协议跟踪：1 入队后，3 进入时弹出 1，因为 3 更大且更晚；窗口形成后队首 3 给出最大值。

\textbf{代码落点。} 关键是队列保存下标而不是值，这样才能判断过期。顺序通常是清队尾、入队、清队首过期、读答案。关键代码行必须能逐句翻译成“旧状态怎样变成新状态”，否则就只是模板。

\textbf{正确性。} 被队尾新元素弹出的旧元素既更小又更早，在未来任何包含旧元素的窗口里也会包含新元素，因此不会成为最大值。

\textbf{边界实验。} 先用k 为一、k 等于数组长度、重复最大值、队首过期做反例检查。实验时准备k 为一、k 等于数组长度、重复最大值、队首过期，打印每次入栈、弹栈和结算对象。重点观察相等元素、空栈、哨兵和最后一次结算。

\textbf{复杂度变化。} 暴力 O(nk)，单调队列每个下标入队出队一次，O(n) 时间、O(k) 空间。

\textbf{迁移追问。} 如果题目改成“受限子序列和使用同样队列维护最近 k 个 DP 最大值”，先判断原状态是否仍然覆盖新答案集合，再决定是微调边界、改 value 语义，还是换算法。

\subsection{LeetCode 49 字母异位词分组}

\textbf{题目说明。} LeetCode 49 字母异位词分组 的题面入口要先还原成具体任务：异位词共享同一个规范表示，可以是排序字符串或字符计数。

\textbf{样例入口。} 拿 eat、tea、tan 手算，eat 和 tea 排序后都是 aet，或计数向量都是 a1e1t1；tan 的键不同。两两比较会重复比较很多字符串，哈希分桶只比较规范键。

\textbf{暴力基线。} 暴力两两比较字符串是否互为异位词，或者对每个新字符串扫描已有分组找能放入的组。

\textbf{暴力过程。} 以 eat、tea、tan 为例，eat 和 tea 排序后都是 aet，应放一组；tan 排序后是 ant，属于另一组。

\textbf{重复工作。} 题面模型：异位词共享同一个规范表示，可以是排序字符串或字符计数。暴力版的慢点要落到本题：浪费在重复比较字符串。异位词的本质是字符频次相同，只要把每个字符串转成规范 key，就能一次定位分组。后面的优化只消掉这类重复工作，不能改变答案集合。

\textbf{优化条件。} 本题的关键观察是：哈希表按规范键分桶，避免两两比较字符串。这句话必须能翻译成可维护状态；如果题目条件改变后观察不再成立，就不能继续套原来的写法。

\textbf{相邻状态。} 规律是异位词共享规范表示，字符集固定时计数键更直接；空字符串和重复单词也必须落到同一个键语义里。把这个特例继续拆开看：每个合法答案都要还原成当前状态和某个历史状态的配对；哈希表的 key 负责定位历史状态，value 负责表达次数、最早位置或存在性。先查还是先写，必须由是否允许当前前缀和自己配对来决定。

\textbf{状态复用。} 哈希表 group[key] 保存同一个规范 key 下的所有原字符串。key 可以是排序后的字符串，也可以是 26 维计数向量编码。手算时只改被新输入影响的那一小部分，而不是回头重算完整候选。

\textbf{指针和字段。} 状态字段要能说清：word 是当前字符串，key 是它的规范表示，groups[key] 是该异位词类的结果数组。手算时按这个协议跟踪：扫描 eat 时创建 key=aet 的组；扫描 tea 时得到同一个 key，追加进同组；扫描 tan 时进入 key=ant 的新组。

\textbf{代码落点。} 排序 key 写法简单但每个词 O(L log L)；计数 key 是 O(L)，但编码时必须加分隔符，避免 1 和 11 这类计数串歧义。关键代码行必须能逐句翻译成“旧状态怎样变成新状态”，否则就只是模板。

\textbf{正确性。} 两个字符串互为异位词当且仅当字符频次完全相同；规范 key 完整表达频次，所以同 key 等价、不同 key 不等价。

\textbf{边界实验。} 先用空字符串、重复单词、字符集不止小写字母、计数键必须无歧义做反例检查。实验时把空字符串、重复单词、字符集不止小写字母、计数键必须无歧义加入对拍集合，用平方暴力和优化版比较。失败时先看初始历史状态、查询更新顺序和哈希表 value 类型。

\textbf{复杂度变化。} 两两比较最坏 O(\ensuremath{n^{2}}L)；排序 key 为 O(nL log L)，计数 key 为 O(nL)，空间 O(nL)。

\textbf{迁移追问。} 如果题目改成“若字符串很长且字符集固定，计数键通常比排序键更稳定”，先判断原状态是否仍然覆盖新答案集合，再决定是微调边界、改 value 语义，还是换算法。

\subsection{LeetCode 128 最长连续序列}

\textbf{题目说明。} LeetCode 128 最长连续序列 的题面入口要先还原成具体任务：连续段只需要从没有前驱的数开始扩展。

\textbf{样例入口。} 拿 [100, 4, 200, 1, 3, 2] 手算，4 不是起点，因为 3 存在；1 没有前驱，从 1 扩到 4 得到长度 4。

\textbf{暴力基线。} 暴力对每个数 x，向上查找 x+1、x+2 是否存在，计算以 x 开头的连续长度。

\textbf{暴力过程。} 以 [100, 4, 200, 1, 3, 2] 为例，若从 2、3、4 都向后扩，会反复扫描同一段 1-4。

\textbf{重复工作。} 题面模型：连续段只需要从没有前驱的数开始扩展。暴力版的慢点要落到本题：浪费在每个连续段被多个内部元素重复扩展。一个连续段只应该从没有前驱 x-1 的起点开始扩。后面的优化只消掉这类重复工作，不能改变答案集合。

\textbf{优化条件。} 本题的关键观察是：哈希集合判断 x 减一 是否存在，跳过非起点避免重复扫描。这句话必须能翻译成可维护状态；如果题目条件改变后观察不再成立，就不能继续套原来的写法。

\textbf{相邻状态。} 规律是只从没有 x-1 的数开始扩展，避免每个连续段被重复扫描；重复元素先由集合去掉。把这个特例继续拆开看：每个合法答案都要还原成当前状态和某个历史状态的配对；哈希表的 key 负责定位历史状态，value 负责表达次数、最早位置或存在性。先查还是先写，必须由是否允许当前前缀和自己配对来决定。

\textbf{状态复用。} 哈希集合保存所有数字。遍历每个 x，只有当 x-1 不存在时，才从 x 开始不断找 x+1。手算时只改被新输入影响的那一小部分，而不是回头重算完整候选。

\textbf{指针和字段。} 状态字段要能说清：num 是候选起点，current 是扩展中的数，length 是当前连续段长度，best 是最大长度。手算时按这个协议跟踪：4 有前驱 3，不做起点；1 没有前驱，从 1 扩到 2、3、4，长度 4。

\textbf{代码落点。} 关键是 if (!set.contains(x-1)) 才进入 while。重复元素先被集合去掉，不影响长度。关键代码行必须能逐句翻译成“旧状态怎样变成新状态”，否则就只是模板。

\textbf{正确性。} 每个连续段有且只有一个没有前驱的最小元素；算法只从这个元素扩展，覆盖每段且不重复扩展内部点。

\textbf{边界实验。} 先用重复元素、负数、空数组、整数边界做反例检查。实验时把重复元素、负数、空数组、整数边界加入对拍集合，用平方暴力和优化版比较。失败时先看初始历史状态、查询更新顺序和哈希表 value 类型。

\textbf{复杂度变化。} 排序法 O(n log n)；哈希起点法均摊 O(n) 时间、O(n) 空间。

\textbf{迁移追问。} 如果题目改成“若数据流动态插入，可以用并查集或区间合并维护长度”，先判断原状态是否仍然覆盖新答案集合，再决定是微调边界、改 value 语义，还是换算法。

\subsection{LeetCode 560 和为 K 的子数组}

\textbf{题目说明。} LeetCode 560 和为 K 的子数组给定整数数组 nums 和整数 k，统计和等于 k 的连续子数组个数；数组中可以有负数。

\textbf{样例入口。} 样例 nums=[1, 1, 1]，k=2，输出 2；两个合法连续子数组分别是前两个 1 和后两个 1。

\textbf{暴力基线。} 暴力固定 left，向右累计 sum；每当 sum==k 就计数。

\textbf{暴力过程。} 以 nums=[1, 1, 1]、k=2 为例，从 left=0 可找到 [1, 1]，从 left=1 又找到后两个 1。

\textbf{重复工作。} 题面模型：当前前缀和减去目标值，就是需要匹配的历史前缀。暴力版的慢点要落到本题：浪费在每个右端都回头枚举所有 left。区间和等于 k 可以改写成 prefix[right]-prefix[left]=k。后面的优化只消掉这类重复工作，不能改变答案集合。

\textbf{优化条件。} 本题的关键观察是：哈希表保存前缀和出现次数，先查询再更新。这句话必须能翻译成可维护状态；如果题目条件改变后观察不再成立，就不能继续套原来的写法。

\textbf{相邻状态。} 规律是哈希表保存前缀和频次，且初始要放 0:1。把这个特例继续拆开看：每个合法答案都要还原成当前状态和某个历史状态的配对；哈希表的 key 负责定位历史状态，value 负责表达次数、最早位置或存在性。先查还是先写，必须由是否允许当前前缀和自己配对来决定。

\textbf{状态复用。} 扫描前缀和 prefix。哈希表 count[old\_prefix] 保存历史前缀出现次数；当前位置需要查 prefix-k 出现过几次。手算时只改被新输入影响的那一小部分，而不是回头重算完整候选。

\textbf{指针和字段。} 状态字段要能说清：prefix 是当前前缀和，need=prefix-k 是所需历史前缀，count 保存历史前缀频次，answer 累加匹配次数。手算时按这个协议跟踪：[1, 1, 1] 中扫描到第二个元素时 prefix=2，需要历史 0，count[0]=1，所以找到第一个长度 2 子数组。

\textbf{代码落点。} 关键是先 answer += count[prefix-k]，再 count[prefix]++。初始化 count[0]=1，表示从下标 0 开始的区间。关键代码行必须能逐句翻译成“旧状态怎样变成新状态”，否则就只是模板。

\textbf{正确性。} 任意和为 k 的子数组都对应一对前缀差 k；扫描到右端时，所有合法左端前缀都已在表里。

\textbf{边界实验。} 先用负数、目标为零、从下标零开始的区间、重复前缀做反例检查。实验时把负数、目标为零、从下标零开始的区间、重复前缀加入对拍集合，用平方暴力和优化版比较。失败时先看初始历史状态、查询更新顺序和哈希表 value 类型。

\textbf{复杂度变化。} 暴力 O(\ensuremath{n^{2}})，前缀哈希 O(n) 均摊时间、O(n) 空间。

\textbf{迁移追问。} 如果题目改成“若要求最长区间，哈希表应保存最早位置而不是次数”，先判断原状态是否仍然覆盖新答案集合，再决定是微调边界、改 value 语义，还是换算法。

\subsection{LeetCode 347 前 K 个高频元素}

\textbf{题目说明。} LeetCode 347 前 K 个高频元素 的题面入口要先还原成具体任务：先统计频次，再只维护前 K 个候选。

\textbf{样例入口。} 拿 [1, 1, 1, 2, 2, 3]、k=2 手算，频次是 1->3、2->2、3->1，前两个高频是 1 和 2。

\textbf{暴力基线。} 暴力先统计频次，再把所有不同元素按频次完整排序，取前 K 个。

\textbf{暴力过程。} 以 [1, 1, 1, 2, 2, 3]、k=2 为例，频次为 1->3、2->2、3->1，答案是 1 和 2。

\textbf{重复工作。} 题面模型：先统计频次，再只维护前 K 个候选。暴力版的慢点要落到本题：浪费在完整排序所有不同元素。只要前 K 高频，低于当前第 K 高频边界的元素可以丢弃。后面的优化只消掉这类重复工作，不能改变答案集合。

\textbf{优化条件。} 本题的关键观察是：小顶堆大小超过 K 就弹出最低频；桶排序利用频次上界。这句话必须能翻译成可维护状态；如果题目条件改变后观察不再成立，就不能继续套原来的写法。

\textbf{相邻状态。} 规律是先用哈希计数，再用大小为 k 的小顶堆或桶按频次取 Top K；频次相同不影响题目集合答案。把这个特例继续拆开看：堆顶必须正好代表下一步最需要处理的候选；若堆里可能有过期对象，读取堆顶前要先清理。堆只解决动态极值，不自动证明被弹出或被丢弃的元素无用。

\textbf{状态复用。} 先用哈希表计数。再用大小为 K 的小顶堆保存候选，堆顶是当前前 K 中频次最低者；超过 K 就弹出。手算时只改被新输入影响的那一小部分，而不是回头重算完整候选。

\textbf{指针和字段。} 状态字段要能说清：freq[value] 是频次，heap 元素是 (frequency, value)，heap.top 是淘汰边界。手算时按这个协议跟踪：插入频次 3 和 2 后堆满；频次 1 的元素入堆会被立刻弹出，留下 1 和 2。

\textbf{代码落点。} 关键是比较器按频次建小顶堆。题目通常不要求输出顺序；若要求顺序，最后还要按规则排序结果。关键代码行必须能逐句翻译成“旧状态怎样变成新状态”，否则就只是模板。

\textbf{正确性。} 堆始终保存已处理元素中频次最高的 K 个；任何被弹出的元素频次不高于堆中 K 个候选，不可能进入答案。

\textbf{边界实验。} 先用K 等于不同元素数、频次相同、结果顺序不要求做反例检查。实验时覆盖K 等于不同元素数、频次相同、结果顺序不要求，记录堆顶是否过期、比较器是否让正确候选在堆顶、K 或窗口大小为边界值时是否稳定。

\textbf{复杂度变化。} 完整排序 O(m log m)，m 为不同元素数；堆法 O(n+m log k)，空间 O(m+k)。

\textbf{迁移追问。} 如果题目改成“若要按频次和字典序排序，比较器要同时处理两个键”，先判断原状态是否仍然覆盖新答案集合，再决定是微调边界、改 value 语义，还是换算法。

\subsection{LeetCode 215 数组中的第 K 个最大元素}

\textbf{题目说明。} LeetCode 215 数组中的第 K 个最大元素 的题面入口要先还原成具体任务：只关心第 K 大，不需要完整排序。

\textbf{样例入口。} 拿 [3, 2, 1, 5, 6, 4]、k=2 手算，固定大小为 2 的小顶堆最终保存最大的两个数，堆顶就是第 2 大。

\textbf{暴力基线。} 暴力排序整个数组，然后返回降序第 k 个或升序第 n-k 个。

\textbf{暴力过程。} 以 [3, 2, 1, 5, 6, 4]、k=2 为例，完整排序后第 2 大是 5；但为了一个第 K 大，不需要知道所有元素的完整顺序。

\textbf{重复工作。} 题面模型：只关心第 K 大，不需要完整排序。暴力版的慢点要落到本题：浪费在完整排序。我们只需要维护当前最大的 K 个元素，其中最小的那个就是第 K 大边界。后面的优化只消掉这类重复工作，不能改变答案集合。

\textbf{优化条件。} 本题的关键观察是：大小为 K 的小顶堆保存当前前 K 大边界；快速选择可做到平均线性。这句话必须能翻译成可维护状态；如果题目条件改变后观察不再成立，就不能继续套原来的写法。

\textbf{相邻状态。} 规律是堆里只保留当前前 k 大候选，比堆顶小的数不可能进入答案。把这个特例继续拆开看：堆顶必须正好代表下一步最需要处理的候选；若堆里可能有过期对象，读取堆顶前要先清理。堆只解决动态极值，不自动证明被弹出或被丢弃的元素无用。

\textbf{状态复用。} 大小为 K 的小顶堆保存当前前 K 大。新数大于堆顶时，弹出堆顶并加入新数；小于等于堆顶则不可能进入前 K。手算时只改被新输入影响的那一小部分，而不是回头重算完整候选。

\textbf{指针和字段。} 状态字段要能说清：heap.top 是当前前 K 大中的最小值，size 保持不超过 K，num 是当前扫描元素。手算时按这个协议跟踪：扫描样例时，堆最终保存 5 和 6，堆顶 5 就是第 2 大。

\textbf{代码落点。} 关键是使用小顶堆，C++ priority\_queue 默认大顶堆，需要 greater<int>。若选择快速选择，要处理随机 pivot 和重复值。关键代码行必须能逐句翻译成“旧状态怎样变成新状态”，否则就只是模板。

\textbf{正确性。} 堆中始终保存已扫描元素的前 K 大；若新元素不超过堆顶，它最多排在第 K 大之后，可以安全丢弃。

\textbf{边界实验。} 先用K 为一、K 为 n、重复值、随机 pivot做反例检查。实验时覆盖K 为一、K 为 n、重复值、随机 pivot，记录堆顶是否过期、比较器是否让正确候选在堆顶、K 或窗口大小为边界值时是否稳定。

\textbf{复杂度变化。} 排序 O(n log n)，小顶堆 O(n log k) 时间、O(k) 空间；快速选择平均 O(n)。

\textbf{迁移追问。} 如果题目改成“若数据流持续到来，堆方案比一次性快速选择更自然”，先判断原状态是否仍然覆盖新答案集合，再决定是微调边界、改 value 语义，还是换算法。

\subsection{LeetCode 102 二叉树层序遍历}

\textbf{题目说明。} LeetCode 102 二叉树层序遍历 的题面入口要先还原成具体任务：队列在每轮开始时保存当前层所有节点。

\textbf{样例入口。} 拿根 3、左 9、右 20 手算，队列第一轮只有 3，第二轮才是 9 和 20。若不固定本层 size，新入队的下一层节点会混进当前层。

\textbf{暴力基线。} 慢版本可以先给每个节点计算深度，再按深度分组；或者每层都从根重新扫描一遍收集该层节点。

\textbf{暴力过程。} 以根 3、左右孩子 9 和 20 为例，第一层只有 3，第二层是 9、20。若队列不固定当前层长度，20 的孩子可能混进同一层。

\textbf{重复工作。} 题面模型：队列在每轮开始时保存当前层所有节点。暴力版的慢点要落到本题：浪费在重复从根找层，或者层边界不清导致混层。队列天然按访问顺序保存下一批节点，只要固定本轮队列大小就是当前层。后面的优化只消掉这类重复工作，不能改变答案集合。

\textbf{优化条件。} 本题的关键观察是：记录 level size 固定层边界，避免下一层节点混入当前层。这句话必须能翻译成可维护状态；如果题目条件改变后观察不再成立，就不能继续套原来的写法。

\textbf{相邻状态。} 规律是每轮开始记录队列长度，这个长度就是当前层边界。把这个特例继续拆开看：节点向父亲返回的量和全局答案通常不是同一个量。先写叶子或空节点的返回值，再写父节点如何合并左右孩子，递归式才有边界基础。

\textbf{状态复用。} BFS 队列保存待处理节点。每轮先读取 level\_size，只弹出这 level\_size 个节点并收集值，同时把它们的孩子加入队尾。手算时只改被新输入影响的那一小部分，而不是回头重算完整候选。

\textbf{指针和字段。} 状态字段要能说清：queue 是待访问节点队列，level\_size 是当前层节点数，level\_values 是本层输出。手算时按这个协议跟踪：队列初始 [3]，level\_size=1，弹出 3 后压入 9 和 20；下一轮 level\_size=2，只处理 9 和 20。

\textbf{代码落点。} 关键是 for 循环次数使用进入本层前的 queue.size()，不要在循环条件里直接用动态变化的 queue.size()。关键代码行必须能逐句翻译成“旧状态怎样变成新状态”，否则就只是模板。

\textbf{正确性。} BFS 入队顺序保证父层节点先于子层节点；固定 level\_size 后，本轮处理的恰好是同一深度的全部节点。

\textbf{边界实验。} 先用空树、只有根、最后一层不满、左右顺序做反例检查。实验时覆盖空树、只有根、最后一层不满、左右顺序，尤其关注空节点、单节点、链式树和全负值。递归版本和显式栈版本可以互相对拍。

\textbf{复杂度变化。} 重复扫描每层最坏 O(nh)，BFS 每节点入队出队一次 O(n)，队列空间 O(width)。

\textbf{迁移追问。} 如果题目改成“锯齿形层序只是在每层输出顺序上增加方向控制”，先判断原状态是否仍然覆盖新答案集合，再决定是微调边界、改 value 语义，还是换算法。

\subsection{LeetCode 98 验证二叉搜索树}

\textbf{题目说明。} LeetCode 98 验证二叉搜索树 的题面入口要先还原成具体任务：BST 约束是全局上下界，不是只比较父子节点。

\textbf{样例入口。} 拿 [5, 1, 4, null, null, 3, 6] 手算，节点 3 小于根 5 却在右子树里，只比较父子 4 和 3 会误判。

\textbf{暴力基线。} 错误暴力常只检查当前节点大于左孩子、小于右孩子；正确慢版本要对每个节点扫描整棵左子树最大值和右子树最小值。

\textbf{暴力过程。} 以 [5, 1, 4, null, null, 3, 6] 为例，节点 4 的左孩子 3 小于 4、右孩子 6 大于 4，看局部没问题；但 3 在根 5 的右子树里，小于 5，所以整棵树不是 BST。

\textbf{重复工作。} 题面模型：BST 约束是全局上下界，不是只比较父子节点。暴力版的慢点要落到本题：浪费在每个节点重复扫描子树求上下界。真正需要传递的是祖先给当前子树限定的取值范围。后面的优化只消掉这类重复工作，不能改变答案集合。

\textbf{优化条件。} 本题的关键观察是：中序严格递增或显式传递上下界都可以；中序版本要保存前一个访问值。这句话必须能翻译成可维护状态；如果题目条件改变后观察不再成立，就不能继续套原来的写法。

\textbf{相邻状态。} 规律是 BST 约束来自整条祖先上下界，或中序序列必须严格递增；重复值和 int 边界要按题意单独处理。把这个特例继续拆开看：节点向父亲返回的量和全局答案通常不是同一个量。先写叶子或空节点的返回值，再写父节点如何合并左右孩子，递归式才有边界基础。

\textbf{状态复用。} 自顶向下传入 lower 和 upper，当前节点值必须满足 lower<val<upper；去左子树时 upper 变成当前值，去右子树时 lower 变成当前值。手算时只改被新输入影响的那一小部分，而不是回头重算完整候选。

\textbf{指针和字段。} 状态字段要能说清：node 是当前节点，lower/upper 是祖先约束；空节点返回合法。手算时按这个协议跟踪：到样例中右子树节点 3 时，它的 lower 应该是 5、upper 是 4 或来自路径约束，3 不满足根给的下界。

\textbf{代码落点。} 关键是用 long long 或 optional 边界，避免节点值等于 int 最小/最大时哨兵冲突；重复值是否允许按题意处理，本题严格不允许。关键代码行必须能逐句翻译成“旧状态怎样变成新状态”，否则就只是模板。

\textbf{正确性。} BST 要求左子树所有值小于根、右子树所有值大于根；上下界把所有祖先约束都带到当前节点，覆盖了非局部违规。

\textbf{边界实验。} 先用重复值、int 最小最大值、局部合法但全局非法做反例检查。实验时覆盖重复值、int 最小最大值、局部合法但全局非法，尤其关注空节点、单节点、链式树和全负值。递归版本和显式栈版本可以互相对拍。

\textbf{复杂度变化。} 扫描子树慢版本最坏 O(\ensuremath{n^{2}})，上下界 DFS/BFS 每节点一次 O(n)，空间 O(h) 或显式栈 O(h)。

\textbf{迁移追问。} 如果题目改成“若允许重复值，需要明确重复放左还是放右”，先判断原状态是否仍然覆盖新答案集合，再决定是微调边界、改 value 语义，还是换算法。

\subsection{LeetCode 236 二叉树最近公共祖先}

\textbf{题目说明。} LeetCode 236 二叉树最近公共祖先 的题面入口要先还原成具体任务：普通树没有有序性，只能依靠祖先关系。

\textbf{样例入口。} 拿普通二叉树根 3，p=5、q=1 手算，左右子树分别找到一个目标，根 3 就是最近公共祖先；若两个目标都在左子树，答案应由左子树返回。

\textbf{暴力基线。} 慢版本在每个节点重新扫描子树或路径，先求出局部答案。重复扫描暴露了真正状态：每个节点只需要从孩子或父亲那里拿到少量信息。放到本题就是：先按“普通树没有有序性，只能依靠祖先关系”完整试慢版本；样例里已经能看到“规律是后序返回是否找到目标或祖先，不能用 BST 的大小关系”，所以优化前必须先能说清慢版本枚举了哪些候选。

\textbf{暴力过程。} 拿普通二叉树根 3，p=5、q=1 手算，左右子树分别找到一个目标，根 3 就是最近公共祖先；若两个目标都在左子树，答案应由左子树返回。

\textbf{重复工作。} 题面模型：普通树没有有序性，只能依靠祖先关系。暴力版的慢点要落到本题：慢版本会在树上重复确认“普通树没有有序性，只能依靠祖先关系”。样例规律是“规律是后序返回是否找到目标或祖先，不能用 BST 的大小关系”，说明父子之间真正需要传递的是高度、上下界、命中状态、层号或两类选择值，而不是反复扫描整棵子树。后面的优化只消掉这类重复工作，不能改变答案集合。

\textbf{优化条件。} 本题的关键观察是：父指针集合或后序返回目标命中状态都可行。这句话必须能翻译成可维护状态；如果题目条件改变后观察不再成立，就不能继续套原来的写法。

\textbf{相邻状态。} 规律是后序返回是否找到目标或祖先，不能用 BST 的大小关系。把这个特例继续拆开看：节点向父亲返回的量和全局答案通常不是同一个量。先写叶子或空节点的返回值，再写父节点如何合并左右孩子，递归式才有边界基础。

\textbf{状态复用。} 把对子树的重复扫描压成返回值或队列状态；每个节点只合并孩子给出的稳定信息。本题落点是：规律是后序返回是否找到目标或祖先，不能用 BST 的大小关系。手算时只改被新输入影响的那一小部分，而不是回头重算完整候选。

\textbf{指针和字段。} 状态字段要能说清：节点指针、传入约束或返回值、当前答案、层号或路径状态；空节点的返回值必须和状态定义匹配。手算时按这个协议跟踪：拿普通二叉树根 3，p=5、q=1 手算，左右子树分别找到一个目标，根 3 就是最近公共祖先；若两个目标都在左子树，答案应由左子树返回。然后按协议复查：手算时从叶子开始写返回值，或者从根开始写传入约束。不要只写访问顺序，要写每条边上传递的信息。

\textbf{代码落点。} 递归版本先处理空节点；迭代版本的栈元素要带完整状态；全负值、重复值和极端深树要单独测。关键代码行必须能逐句翻译成“旧状态怎样变成新状态”，否则就只是模板。

\textbf{正确性。} 证明从子树归纳或层序不变量开始：孩子状态正确时父节点合并正确；若用队列，当前轮队列必须正好保存当前层节点。

\textbf{边界实验。} 先用目标不存在、p 等于 q、一个是另一个祖先、比较节点地址做反例检查。实验时覆盖目标不存在、p 等于 q、一个是另一个祖先、比较节点地址，尤其关注空节点、单节点、链式树和全负值。递归版本和显式栈版本可以互相对拍。

\textbf{复杂度变化。} 暴力在每个节点重复扫描子树或反复寻找层边界可能退化到平方级；后序汇总、显式栈或层序遍历让每个节点处理常数次，通常是线性时间。

\textbf{迁移追问。} 如果题目改成“多次 LCA 查询可预处理倍增祖先表”，先判断原状态是否仍然覆盖新答案集合，再决定是微调边界、改 value 语义，还是换算法。

\subsection{LeetCode 200 岛屿数量}

\textbf{题目说明。} LeetCode 200 岛屿数量 的题面入口要先还原成具体任务：陆地格子是节点，四方向相邻是边，岛屿是连通块。

\textbf{样例入口。} 拿一个 2x2 小岛手算，从一个陆地出发能沿上下左右标记同一块岛。若不标记访问，会重复计数；若对角线也连通，会误合并。

\textbf{暴力基线。} 暴力可以对每个陆地格子单独搜索它所属的连通块，再用集合去重；这样同一块岛会被重复搜索很多次。

\textbf{暴力过程。} 以 2x2 全是陆地为例，从左上搜索会访问四个格子；如果之后从右上又重新搜索，就把同一座岛重复数了一遍。

\textbf{重复工作。} 题面模型：陆地格子是节点，四方向相邻是边，岛屿是连通块。暴力版的慢点要落到本题：浪费在已确认属于某个岛的陆地没有被标记。每块岛只需要在第一次遇到时完整淹没或标记。后面的优化只消掉这类重复工作，不能改变答案集合。

\textbf{优化条件。} 本题的关键观察是：扫描遇到未访问陆地就启动搜索并标记整块。这句话必须能翻译成可维护状态；如果题目条件改变后观察不再成立，就不能继续套原来的写法。

\textbf{相邻状态。} 规律是每发现一个未访问陆地，答案加一，并用 DFS/BFS 标记与它四连通的全部陆地。把这个特例继续拆开看：状态节点不一定等于原始位置，还可能包含钥匙、剩余资源、时间或访问集合。visited 只能合并未来能力完全相同的状态，否则会把合法路径剪掉。

\textbf{状态复用。} 扫描网格。遇到未访问陆地，答案加一，然后 BFS/DFS 把与它四方向连通的全部陆地标记为已访问。手算时只改被新输入影响的那一小部分，而不是回头重算完整候选。

\textbf{指针和字段。} 状态字段要能说清：row/col 是当前格子，visited 或原地改写表示已经归属某座岛，directions 是四方向邻居。手算时按这个协议跟踪：从左上陆地入队，依次扩展右边和下边陆地，最终整块连通区域都标记；扫描到这些格子时直接跳过。

\textbf{代码落点。} 关键是只看上下左右，不看对角线。边界判断集中写，访问标记在入队时设置，避免重复入队。关键代码行必须能逐句翻译成“旧状态怎样变成新状态”，否则就只是模板。

\textbf{正确性。} 一次搜索会访问且只访问与起点四连通的全部陆地；外层每次启动搜索都对应一块此前未计数的新岛。

\textbf{边界实验。} 先用空网格、全水、全陆、斜对角不连通、是否允许修改输入做反例检查。实验时覆盖空网格、全水、全陆、斜对角不连通、是否允许修改输入，并打印入队时的完整状态。若同一位置但资源不同的状态被合并，很多图题会出现隐蔽错误。

\textbf{复杂度变化。} 重复搜索最坏 O(\ensuremath{(mn)^{2}})，标记搜索 O(mn) 时间，空间 O(mn) 或递归/队列大小。

\textbf{迁移追问。} 如果题目改成“也可用并查集合并相邻陆地，适合动态岛屿扩展”，先判断原状态是否仍然覆盖新答案集合，再决定是微调边界、改 value 语义，还是换算法。

\subsection{LeetCode 994 腐烂的橘子}

\textbf{题目说明。} LeetCode 994 腐烂的橘子 的题面入口要先还原成具体任务：所有初始腐烂橘子同时作为 BFS 源点。

\textbf{样例入口。} 拿 [[2, 1, 1], [1, 1, 0], [0, 1, 1]] 手算，所有腐烂橘子同时向四周扩散，第一分钟会感染相邻新鲜橘子。

\textbf{暴力基线。} 暴力按分钟模拟整张网格：每一分钟扫描所有格子，找到会被感染的新鲜橘子并统一变腐烂。

\textbf{暴力过程。} 在 [[2, 1, 1], [1, 1, 0], [0, 1, 1]] 中，初始腐烂橘子同时向四周扩散；第一分钟感染它上下左右的新鲜橘子。

\textbf{重复工作。} 题面模型：所有初始腐烂橘子同时作为 BFS 源点。暴力版的慢点要落到本题：浪费在每分钟都扫描全图。扩散只会从上一分钟刚腐烂的边界继续向外，队列正好保存这个边界。后面的优化只消掉这类重复工作，不能改变答案集合。

\textbf{优化条件。} 本题的关键观察是：按层扩散，分钟数等于最后一个新鲜橘子被首次访问的层数。这句话必须能翻译成可维护状态；如果题目条件改变后观察不再成立，就不能继续套原来的写法。

\textbf{相邻状态。} 规律是多源 BFS，初始所有腐烂橘子同时入队，层数就是分钟数。把这个特例继续拆开看：状态节点不一定等于原始位置，还可能包含钥匙、剩余资源、时间或访问集合。visited 只能合并未来能力完全相同的状态，否则会把合法路径剪掉。

\textbf{状态复用。} 多源 BFS。初始把所有腐烂橘子入队，并统计 fresh 数量；每弹出一层代表一分钟，感染相邻新鲜橘子并 fresh--。手算时只改被新输入影响的那一小部分，而不是回头重算完整候选。

\textbf{指针和字段。} 状态字段要能说清：queue 保存当前扩散边界，fresh 是剩余新鲜橘子数，minutes 是已经完成的层数，directions 是四方向。手算时按这个协议跟踪：第一层处理初始腐烂源，感染邻居入队；第二层只处理第一分钟新腐烂的橘子，层数自然对应分钟。

\textbf{代码落点。} 关键是所有初始腐烂橘子一次性入队。感染时立即把 grid[nr][nc] 改成 2，避免同一分钟被重复入队。关键代码行必须能逐句翻译成“旧状态怎样变成新状态”，否则就只是模板。

\textbf{正确性。} BFS 按层扩展，某个新鲜橘子第一次被访问的层数就是它能被腐烂的最早分钟；多源初始化表示所有源同时开始。

\textbf{边界实验。} 先用没有新鲜橘子、没有腐烂源、隔离新鲜橘子、空格阻断做反例检查。实验时覆盖没有新鲜橘子、没有腐烂源、隔离新鲜橘子、空格阻断，并打印入队时的完整状态。若同一位置但资源不同的状态被合并，很多图题会出现隐蔽错误。

\textbf{复杂度变化。} 逐分钟全图扫描最坏 O(mn*分钟数)，BFS 每格最多入队一次，O(mn) 时间、O(mn) 空间。

\textbf{迁移追问。} 如果题目改成“多源 BFS 也适用于最近零距离、地图扩散等题”，先判断原状态是否仍然覆盖新答案集合，再决定是微调边界、改 value 语义，还是换算法。

\subsection{LeetCode 207 课程表}

\textbf{题目说明。} LeetCode 207 课程表 的题面入口要先还原成具体任务：课程是节点，先修关系是有向边，问题是判断是否有环。

\textbf{样例入口。} 拿课程 0->1 手算，入度为 0 的课程 0 先入队，学完后 1 的入度降为 0；若有 0->1 和 1->0，队列一开始为空。

\textbf{暴力基线。} 暴力可以不断扫描所有课程，找当前先修都已满足的课程；若一轮没有新课程可学，就判断有环。

\textbf{暴力过程。} 以 0->1 为例，课程 0 入度为 0，先学 0，随后 1 的入度变 0。若有 0->1 和 1->0，一开始没有入度为 0 的课。

\textbf{重复工作。} 题面模型：课程是节点，先修关系是有向边，问题是判断是否有环。暴力版的慢点要落到本题：浪费在每轮都重新检查所有课程的先修。入度正好表示还剩多少未满足依赖，学完一门课只影响它指向的后继。后面的优化只消掉这类重复工作，不能改变答案集合。

\textbf{优化条件。} 本题的关键观察是：入度拓扑排序不断学习当前无依赖课程。这句话必须能翻译成可维护状态；如果题目条件改变后观察不再成立，就不能继续套原来的写法。

\textbf{相邻状态。} 规律是拓扑排序用入度为 0 的节点推进，处理数量小于课程数说明有环。把这个特例继续拆开看：状态节点不一定等于原始位置，还可能包含钥匙、剩余资源、时间或访问集合。visited 只能合并未来能力完全相同的状态，否则会把合法路径剪掉。

\textbf{状态复用。} 构建邻接表和 indegree。把所有入度为 0 的课程入队，弹出课程时让后继入度减一，减到 0 就入队。手算时只改被新输入影响的那一小部分，而不是回头重算完整候选。

\textbf{指针和字段。} 状态字段要能说清：indegree[v] 是课程 v 尚未满足的先修数量，queue 保存当前可学习课程，visited\_count 是已处理数量。手算时按这个协议跟踪：0 出队后，遍历它的后继 1，indegree[1] 从 1 变 0，1 入队；最后处理数量等于课程数，说明无环。

\textbf{代码落点。} 关键是边方向要统一：prerequisite [a, b] 表示 b -> a。处理数量小于 numCourses 时返回 false。关键代码行必须能逐句翻译成“旧状态怎样变成新状态”，否则就只是模板。

\textbf{正确性。} 入度为 0 的节点没有未满足前置条件，可以安全学习；有向环中的节点永远不会全部降到入度 0。

\textbf{边界实验。} 先用边方向、孤立课程、环、自依赖、重复边做反例检查。实验时覆盖边方向、孤立课程、环、自依赖、重复边，并打印入队时的完整状态。若同一位置但资源不同的状态被合并，很多图题会出现隐蔽错误。

\textbf{复杂度变化。} 反复扫描最坏 O(VE)，拓扑排序 O(V+E) 时间和空间。

\textbf{迁移追问。} 如果题目改成“DFS 三色标记也能判环，但要明确访问状态”，先判断原状态是否仍然覆盖新答案集合，再决定是微调边界、改 value 语义，还是换算法。

\subsection{LeetCode 210 课程表 II}

\textbf{题目说明。} LeetCode 210 课程表 II 的题面入口要先还原成具体任务：在判环基础上还要输出一个可行拓扑序。

\textbf{样例入口。} 拿 0->1、0->2 手算，课程 0 入度为 0，先输出 0，再把 1 和 2 的入度减到 0。若存在环，输出数量会少于课程数。

\textbf{暴力基线。} 慢版本先按题意画出完整状态图，用普通搜索跑通可达性。这个过程用于确认在判环基础上还要输出一个可行拓扑序的节点和边，之后再讨论访问标记、层次或代价顺序。放到本题就是：先按“在判环基础上还要输出一个可行拓扑序”完整试慢版本；样例里已经能看到“规律是课程表 II 比判定题多了输出顺序，队列弹出的顺序就是一种可行拓扑序”，所以优化前必须先能说清慢版本枚举了哪些候选。

\textbf{暴力过程。} 拿 0->1、0->2 手算，课程 0 入度为 0，先输出 0，再把 1 和 2 的入度减到 0。若存在环，输出数量会少于课程数。

\textbf{重复工作。} 题面模型：在判环基础上还要输出一个可行拓扑序。暴力版的慢点要落到本题：慢版本会围绕“在判环基础上还要输出一个可行拓扑序”反复展开同一批状态。样例规律是“规律是课程表 II 比判定题多了输出顺序，队列弹出的顺序就是一种可行拓扑序”，说明必须把节点、边、访问标记和资源维度一次建清；同一状态不应重复入队，不同资源状态也不能误合并。后面的优化只消掉这类重复工作，不能改变答案集合。

\textbf{优化条件。} 本题的关键观察是：每弹出一个入度为零课程，就把它加入顺序并删除出边。这句话必须能翻译成可维护状态；如果题目条件改变后观察不再成立，就不能继续套原来的写法。

\textbf{相邻状态。} 规律是课程表 II 比判定题多了输出顺序，队列弹出的顺序就是一种可行拓扑序。把这个特例继续拆开看：状态节点不一定等于原始位置，还可能包含钥匙、剩余资源、时间或访问集合。visited 只能合并未来能力完全相同的状态，否则会把合法路径剪掉。

\textbf{状态复用。} 把重复搜索压成访问标记、距离数组或拓扑状态；邻居生成也要从全量扫描变成结构化枚举。本题落点是：规律是课程表 II 比判定题多了输出顺序，队列弹出的顺序就是一种可行拓扑序。手算时只改被新输入影响的那一小部分，而不是回头重算完整候选。

\textbf{指针和字段。} 状态字段要能说清：状态编号、邻接生成方式、访问标记、距离或层数、队列内容；若状态含资源，visited 不能只按位置记录。手算时按这个协议跟踪：拿 0->1、0->2 手算，课程 0 入度为 0，先输出 0，再把 1 和 2 的入度减到 0。若存在环，输出数量会少于课程数。然后按协议复查：手算时画队列或栈，状态必须包含题目需要的所有资源。入队时标记还是出队时标记，要和去重语义一致。

\textbf{代码落点。} 入队时标记还是出队时标记必须一致；方向数组集中定义；多源 BFS 要一次性把所有源点放入初始队列。关键代码行必须能逐句翻译成“旧状态怎样变成新状态”，否则就只是模板。

\textbf{正确性。} 证明从可达性开始：搜索覆盖所有合法边能到达的状态，访问标记只删除重复状态而不删除不同未来能力。

\textbf{边界实验。} 先用有环返回空数组、多个可行顺序、边方向做反例检查。实验时覆盖有环返回空数组、多个可行顺序、边方向，并打印入队时的完整状态。若同一位置但资源不同的状态被合并，很多图题会出现隐蔽错误。

\textbf{复杂度变化。} 暴力重复从多个起点搜索会反复遍历图；正确建模和访问标记后，BFS 或 DFS 通常按节点数加边数计算，状态图题则按完整状态数计算。

\textbf{迁移追问。} 如果题目改成“若需要字典序最小拓扑序，把队列换成小顶堆”，先判断原状态是否仍然覆盖新答案集合，再决定是微调边界、改 value 语义，还是换算法。

\subsection{LeetCode 46 全排列}

\textbf{题目说明。} LeetCode 46 全排列 的题面入口要先还原成具体任务：排列关心元素顺序，每个位置从所有未使用元素中选择一个。

\textbf{样例入口。} 拿 [1, 2, 3] 手算，第一位可选 1、2、3；选 1 后，第二位只能从未使用的 2、3 中选。

\textbf{暴力基线。} 暴力可以枚举所有长度为 n 的序列，再检查是否每个元素刚好出现一次。

\textbf{暴力过程。} 以 [1, 2, 3] 为例，第一位可选 1、2、3；若第一位选 1，第二位只能从 2、3 中选，第三位由剩余元素决定。

\textbf{重复工作。} 题面模型：排列关心元素顺序，每个位置从所有未使用元素中选择一个。暴力版的慢点要落到本题：浪费在生成了含重复使用元素的非法序列。排列的每一层只需要从未使用元素中选择。后面的优化只消掉这类重复工作，不能改变答案集合。

\textbf{优化条件。} 本题的关键观察是：路径长度达到 n 时产生答案，used 数组保证同一元素不重复使用。这句话必须能翻译成可维护状态；如果题目条件改变后观察不再成立，就不能继续套原来的写法。

\textbf{相邻状态。} 规律是路径长度等于 n 时得到一个排列，used 数组表示哪些下标已经在当前路径里。把这个特例继续拆开看：一层递归对应一个明确决策，路径保存已经做出的选择，选择列表保存仍可能扩展的分支。剪枝前必须能说明被剪掉的整棵子树没有合法答案，而不是只因为它看起来不优。

\textbf{状态复用。} 回溯路径 path 保存当前排列前缀，used[i] 表示 nums[i] 是否已在路径中。path 长度等于 n 时复制进答案。手算时只改被新输入影响的那一小部分，而不是回头重算完整候选。

\textbf{指针和字段。} 状态字段要能说清：depth 是当前要填的位置，path 是已选择序列，used 是可选状态，answers 保存所有完整排列。手算时按这个协议跟踪：[1, 2, 3] 中路径 [1] 后，used[1]=true；下一层只能选 2 或 3。回退时弹出 1 并恢复 used。

\textbf{代码落点。} 关键是进入分支前 path.push\_back(nums[i]); used[i]=true；退出分支时反向恢复。保存答案时必须复制 path。关键代码行必须能逐句翻译成“旧状态怎样变成新状态”，否则就只是模板。

\textbf{正确性。} 每个位置都枚举所有未使用元素，所以任意合法排列都有唯一生成路径；used 阻止同一元素在一条路径中重复出现。

\textbf{边界实验。} 先用空数组、单元素、输出规模为 n 阶乘、元素值唯一假设做反例检查。实验时覆盖空数组、单元素、输出规模为 n 阶乘、元素值唯一假设，统计搜索节点数量和剪枝次数。重复元素题要打印同一层跳过哪些值，确认没有误删合法路径。

\textbf{复杂度变化。} 输出规模为 n! 个排列，每个排列长度 n，时间 O(n*n!)，空间 O(n) 不计输出。

\textbf{迁移追问。} 如果题目改成“若有重复元素，需要排序并在同层跳过等价选择”，先判断原状态是否仍然覆盖新答案集合，再决定是微调边界、改 value 语义，还是换算法。

\subsection{LeetCode 78 子集}

\textbf{题目说明。} LeetCode 78 子集 的题面入口要先还原成具体任务：每个元素都有选或不选两种选择，答案是所有路径节点。

\textbf{样例入口。} 拿 [1, 2] 画搜索树，第一层可以不选 1 或选 1；在选 1 的分支里再决定 2，得到 [1] 和 [1, 2]。

\textbf{暴力基线。} 暴力可以枚举长度从 0 到 n 的所有组合，并为每种长度单独搜索。

\textbf{暴力过程。} 以 [1, 2, 3] 为例，空集、[1]、[2]、[3]、[1, 2] 等都是答案；一个元素只有选和不选两种状态。

\textbf{重复工作。} 题面模型：每个元素都有选或不选两种选择，答案是所有路径节点。暴力版的慢点要落到本题：浪费在按长度重复组织搜索。子集问题的每个搜索树节点本身就是一个答案，不必等到固定长度。后面的优化只消掉这类重复工作，不能改变答案集合。

\textbf{优化条件。} 本题的关键观察是：回溯按起点向后枚举，避免不同顺序生成同一集合。这句话必须能翻译成可维护状态；如果题目条件改变后观察不再成立，就不能继续套原来的写法。

\textbf{相邻状态。} 规律是每个元素对应选或不选，或用起点向后枚举，路径上的每个节点都是一个子集；空数组也要输出空集。把这个特例继续拆开看：一层递归对应一个明确决策，路径保存已经做出的选择，选择列表保存仍可能扩展的分支。剪枝前必须能说明被剪掉的整棵子树没有合法答案，而不是只因为它看起来不优。

\textbf{状态复用。} 回溯时先把当前 path 加入答案，再从 start 到末尾选择下一个元素，递归时 start=i+1，保证元素顺序递增。手算时只改被新输入影响的那一小部分，而不是回头重算完整候选。

\textbf{指针和字段。} 状态字段要能说清：start 是下一次可选择的最小下标，path 是当前子集，answers 保存所有已访问路径节点。手算时按这个协议跟踪：从空集开始记录 []；选择 1 后记录 [1]；再选择 2 记录 [1, 2]；回退后选择 3 得到 [1, 3]。

\textbf{代码落点。} 关键是每进入一个节点就记录 path，而不是只在叶子记录。递归下一层用 i+1，避免 [1, 2] 和 [2, 1] 重复。关键代码行必须能逐句翻译成“旧状态怎样变成新状态”，否则就只是模板。

\textbf{正确性。} 任意子集按原数组下标升序排列后，对应搜索树中唯一一条递增路径；每个路径节点都被记录，所以不漏不重。

\textbf{边界实验。} 先用空数组、单元素、输出数量为二的 n 次方、顺序不要求做反例检查。实验时覆盖空数组、单元素、输出数量为二的 n 次方、顺序不要求，统计搜索节点数量和剪枝次数。重复元素题要打印同一层跳过哪些值，确认没有误删合法路径。

\textbf{复杂度变化。} 共有 \ensuremath{2^{n}} 个子集，复制路径总成本 O(\ensuremath{n \cdot 2^{n}})，递归栈 O(n)，输出空间另计。

\textbf{迁移追问。} 如果题目改成“也可用位掩码枚举所有子集，适合 n 较小的场景”，先判断原状态是否仍然覆盖新答案集合，再决定是微调边界、改 value 语义，还是换算法。

\subsection{LeetCode 39 组合总和}

\textbf{题目说明。} LeetCode 39 组合总和 的题面入口要先还原成具体任务：候选数字可重复使用，答案不关心排列顺序。

\textbf{样例入口。} 拿 candidates=[2, 3, 6, 7]、target=7 手算，选择 2 后仍可以继续选择 2，因为每个数可重复；路径 2, 2, 3 成功，7 单独也成功。

\textbf{暴力基线。} 慢版本就是完整搜索树：每层列出选择列表，路径到达终止条件时检查答案。它先保证覆盖所有可能，再把明显无效或重复的分支剪掉。放到本题就是：先按“候选数字可重复使用，答案不关心排列顺序”完整试慢版本；样例里已经能看到“规律是回溯起点不回退，但递归下一层仍传当前下标以允许重复使用”，所以优化前必须先能说清慢版本枚举了哪些候选。

\textbf{暴力过程。} 拿 candidates=[2, 3, 6, 7]、target=7 手算，选择 2 后仍可以继续选择 2，因为每个数可重复；路径 2, 2, 3 成功，7 单独也成功。

\textbf{重复工作。} 题面模型：候选数字可重复使用，答案不关心排列顺序。暴力版的慢点要落到本题：慢版本会围绕“候选数字可重复使用，答案不关心排列顺序”展开完整搜索树。样例规律是“规律是回溯起点不回退，但递归下一层仍传当前下标以允许重复使用”，说明一层递归必须对应一个明确决策，剪枝只能删除整棵不可能成功或重复的子树。后面的优化只消掉这类重复工作，不能改变答案集合。

\textbf{优化条件。} 本题的关键观察是：排序后按起点枚举，递归时仍传当前下标表示可以继续使用同一数字。这句话必须能翻译成可维护状态；如果题目条件改变后观察不再成立，就不能继续套原来的写法。

\textbf{相邻状态。} 规律是回溯起点不回退，但递归下一层仍传当前下标以允许重复使用。把这个特例继续拆开看：一层递归对应一个明确决策，路径保存已经做出的选择，选择列表保存仍可能扩展的分支。剪枝前必须能说明被剪掉的整棵子树没有合法答案，而不是只因为它看起来不优。

\textbf{状态复用。} 把完整搜索树加上合法性检查、剩余资源剪枝和同层去重；路径状态进入和退出要成对恢复。本题落点是：规律是回溯起点不回退，但递归下一层仍传当前下标以允许重复使用。手算时只改被新输入影响的那一小部分，而不是回头重算完整候选。

\textbf{指针和字段。} 状态字段要能说清：路径、起点或使用标记、剩余目标、答案集合、剪枝条件；进入递归和退出递归必须成对恢复。手算时按这个协议跟踪：拿 candidates=[2, 3, 6, 7]、target=7 手算，选择 2 后仍可以继续选择 2，因为每个数可重复；路径 2, 2, 3 成功，7 单独也成功。然后按协议复查：手算时给每层标出选择列表和路径。剪枝发生前要指出被剪分支为什么已经不可能得到合法答案。

\textbf{代码落点。} 剪枝条件写在扩展前；同层去重不要误写成全局去重；保存答案时复制当前路径。关键代码行必须能逐句翻译成“旧状态怎样变成新状态”，否则就只是模板。

\textbf{正确性。} 证明从搜索树覆盖开始：每个合法答案有且只有一条路径，剪枝只删掉不可能成功的分支。

\textbf{边界实验。} 先用目标小于最小数、刚好等于、候选有重复时题意如何处理、输出规模做反例检查。实验时覆盖目标小于最小数、刚好等于、候选有重复时题意如何处理、输出规模，统计搜索节点数量和剪枝次数。重复元素题要打印同一层跳过哪些值，确认没有误删合法路径。

\textbf{复杂度变化。} 暴力搜索树的规模由输出或选择空间决定；剪枝不改变最坏指数级本质，但能删除不可能成功或重复的分支，复杂度要说明输出规模。

\textbf{迁移追问。} 如果题目改成“组合总和 II 每个数只能用一次，起点推进和去重条件都不同”，先判断原状态是否仍然覆盖新答案集合，再决定是微调边界、改 value 语义，还是换算法。

\subsection{LeetCode 22 括号生成}

\textbf{题目说明。} LeetCode 22 括号生成 的题面入口要先还原成具体任务：合法括号串的前缀中右括号数量不能超过左括号数量。

\textbf{样例入口。} 拿 n=2 手算，第一步只能放左括号，之后可以放左括号得到 (())，也可以在已有左括号数量大于右括号时放右括号得到 ()()。

\textbf{暴力基线。} 慢版本就是完整搜索树：每层列出选择列表，路径到达终止条件时检查答案。它先保证覆盖所有可能，再把明显无效或重复的分支剪掉。放到本题就是：先按“合法括号串的前缀中右括号数量不能超过左括号数量”完整试慢版本；样例里已经能看到“规律是状态包含已放左括号数和右括号数，右括号不能超过左括号，左括号不能超过 n”，所以优化前必须先能说清慢版本枚举了哪些候选。

\textbf{暴力过程。} 拿 n=2 手算，第一步只能放左括号，之后可以放左括号得到 (())，也可以在已有左括号数量大于右括号时放右括号得到 ()()。

\textbf{重复工作。} 题面模型：合法括号串的前缀中右括号数量不能超过左括号数量。暴力版的慢点要落到本题：慢版本会围绕“合法括号串的前缀中右括号数量不能超过左括号数量”展开完整搜索树。样例规律是“规律是状态包含已放左括号数和右括号数，右括号不能超过左括号，左括号不能超过 n”，说明一层递归必须对应一个明确决策，剪枝只能删除整棵不可能成功或重复的子树。后面的优化只消掉这类重复工作，不能改变答案集合。

\textbf{优化条件。} 本题的关键观察是：状态保存已用左括号和右括号数量，只有满足约束的选择才继续递归。这句话必须能翻译成可维护状态；如果题目条件改变后观察不再成立，就不能继续套原来的写法。

\textbf{相邻状态。} 规律是状态包含已放左括号数和右括号数，右括号不能超过左括号，左括号不能超过 n。把这个特例继续拆开看：一层递归对应一个明确决策，路径保存已经做出的选择，选择列表保存仍可能扩展的分支。剪枝前必须能说明被剪掉的整棵子树没有合法答案，而不是只因为它看起来不优。

\textbf{状态复用。} 把完整搜索树加上合法性检查、剩余资源剪枝和同层去重；路径状态进入和退出要成对恢复。本题落点是：规律是状态包含已放左括号数和右括号数，右括号不能超过左括号，左括号不能超过 n。手算时只改被新输入影响的那一小部分，而不是回头重算完整候选。

\textbf{指针和字段。} 状态字段要能说清：路径、起点或使用标记、剩余目标、答案集合、剪枝条件；进入递归和退出递归必须成对恢复。手算时按这个协议跟踪：拿 n=2 手算，第一步只能放左括号，之后可以放左括号得到 (())，也可以在已有左括号数量大于右括号时放右括号得到 ()()。然后按协议复查：手算时给每层标出选择列表和路径。剪枝发生前要指出被剪分支为什么已经不可能得到合法答案。

\textbf{代码落点。} 剪枝条件写在扩展前；同层去重不要误写成全局去重；保存答案时复制当前路径。关键代码行必须能逐句翻译成“旧状态怎样变成新状态”，否则就只是模板。

\textbf{正确性。} 证明从搜索树覆盖开始：每个合法答案有且只有一条路径，剪枝只删掉不可能成功的分支。

\textbf{边界实验。} 先用n 为零或一、右括号提前超过左括号、输出规模是卡特兰数做反例检查。实验时覆盖n 为零或一、右括号提前超过左括号、输出规模是卡特兰数，统计搜索节点数量和剪枝次数。重复元素题要打印同一层跳过哪些值，确认没有误删合法路径。

\textbf{复杂度变化。} 暴力搜索树的规模由输出或选择空间决定；剪枝不改变最坏指数级本质，但能删除不可能成功或重复的分支，复杂度要说明输出规模。

\textbf{迁移追问。} 如果题目改成“若有多种括号类型，状态还需要栈保存类型匹配关系”，先判断原状态是否仍然覆盖新答案集合，再决定是微调边界、改 value 语义，还是换算法。

\subsection{LeetCode 79 单词搜索}

\textbf{题目说明。} LeetCode 79 单词搜索 的题面入口要先还原成具体任务：网格路径不能重复使用同一格，状态包含位置、匹配下标和访问标记。

\textbf{样例入口。} 拿 board=[A B; C D]、word=AB 手算，从 A 出发只能走到相邻的 B，不能跳到 D；若搜索 ABA，第二个 A 不能复用起点格。

\textbf{暴力基线。} 慢版本就是完整搜索树：每层列出选择列表，路径到达终止条件时检查答案。它先保证覆盖所有可能，再把明显无效或重复的分支剪掉。放到本题就是：先按“网格路径不能重复使用同一格，状态包含位置、匹配下标和访问标记”完整试慢版本；样例里已经能看到“规律是状态包含位置、匹配下标和访问标记，递归退出时要恢复访问标记”，所以优化前必须先能说清慢版本枚举了哪些候选。

\textbf{暴力过程。} 拿 board=[A B; C D]、word=AB 手算，从 A 出发只能走到相邻的 B，不能跳到 D；若搜索 ABA，第二个 A 不能复用起点格。

\textbf{重复工作。} 题面模型：网格路径不能重复使用同一格，状态包含位置、匹配下标和访问标记。暴力版的慢点要落到本题：慢版本会围绕“网格路径不能重复使用同一格，状态包含位置、匹配下标和访问标记”展开完整搜索树。样例规律是“规律是状态包含位置、匹配下标和访问标记，递归退出时要恢复访问标记”，说明一层递归必须对应一个明确决策，剪枝只能删除整棵不可能成功或重复的子树。后面的优化只消掉这类重复工作，不能改变答案集合。

\textbf{优化条件。} 本题的关键观察是：剪枝包括首字符不匹配、剩余长度不足、字符频次不足和越界访问。这句话必须能翻译成可维护状态；如果题目条件改变后观察不再成立，就不能继续套原来的写法。

\textbf{相邻状态。} 规律是状态包含位置、匹配下标和访问标记，递归退出时要恢复访问标记。把这个特例继续拆开看：一层递归对应一个明确决策，路径保存已经做出的选择，选择列表保存仍可能扩展的分支。剪枝前必须能说明被剪掉的整棵子树没有合法答案，而不是只因为它看起来不优。

\textbf{状态复用。} 把完整搜索树加上合法性检查、剩余资源剪枝和同层去重；路径状态进入和退出要成对恢复。本题落点是：规律是状态包含位置、匹配下标和访问标记，递归退出时要恢复访问标记。手算时只改被新输入影响的那一小部分，而不是回头重算完整候选。

\textbf{指针和字段。} 状态字段要能说清：路径、起点或使用标记、剩余目标、答案集合、剪枝条件；进入递归和退出递归必须成对恢复。手算时按这个协议跟踪：拿 board=[A B; C D]、word=AB 手算，从 A 出发只能走到相邻的 B，不能跳到 D；若搜索 ABA，第二个 A 不能复用起点格。然后按协议复查：手算时给每层标出选择列表和路径。剪枝发生前要指出被剪分支为什么已经不可能得到合法答案。

\textbf{代码落点。} 剪枝条件写在扩展前；同层去重不要误写成全局去重；保存答案时复制当前路径。关键代码行必须能逐句翻译成“旧状态怎样变成新状态”，否则就只是模板。

\textbf{正确性。} 证明从搜索树覆盖开始：每个合法答案有且只有一条路径，剪枝只删掉不可能成功的分支。

\textbf{边界实验。} 先用单字符单词、路径需要回退、重复字母、单元格不能复用做反例检查。实验时覆盖单字符单词、路径需要回退、重复字母、单元格不能复用，统计搜索节点数量和剪枝次数。重复元素题要打印同一层跳过哪些值，确认没有误删合法路径。

\textbf{复杂度变化。} 暴力搜索树的规模由输出或选择空间决定；剪枝不改变最坏指数级本质，但能删除不可能成功或重复的分支，复杂度要说明输出规模。

\textbf{迁移追问。} 如果题目改成“若要查很多单词，应使用 Trie 合并搜索前缀”，先判断原状态是否仍然覆盖新答案集合，再决定是微调边界、改 value 语义，还是换算法。

\subsection{LeetCode 70 爬楼梯}

\textbf{题目说明。} LeetCode 70 爬楼梯 的题面入口要先还原成具体任务：最后一步来自前一阶或前两阶，是最小的重叠子问题模型。

\textbf{样例入口。} 拿 n=4 手算，最后一步要么从第 3 阶跨 1 阶，要么从第 2 阶跨 2 阶，这两类互斥且覆盖所有走法。于是 ways[4]=ways[3]+ways[2]。

\textbf{暴力基线。} 暴力递归枚举每一步走 1 阶或 2 阶，直到正好到 n 计数，超过 n 则失败。

\textbf{暴力过程。} 以 n=4 为例，最后一步要么从第 3 阶走 1 阶，要么从第 2 阶走 2 阶；这两类互斥且覆盖所有走法。

\textbf{重复工作。} 题面模型：最后一步来自前一阶或前两阶，是最小的重叠子问题模型。暴力版的慢点要落到本题：浪费在递归树反复计算 ways(2)、ways(3) 等同一子问题。到达某一阶后，后面的走法数量和之前怎么来无关。后面的优化只消掉这类重复工作，不能改变答案集合。

\textbf{优化条件。} 本题的关键观察是：滚动变量保存最近两个状态即可，不需要完整数组。这句话必须能翻译成可维护状态；如果题目条件改变后观察不再成立，就不能继续套原来的写法。

\textbf{相邻状态。} 规律不是背斐波那契，而是最后一步只有两个来源，状态只需保存前两项。把这个特例继续拆开看：先问暴力搜索里哪些不同路径到达同一个后续问题，再给这个后续问题命名为状态。状态必须包含未来决策需要的全部信息，转移只从更小且已经算清的状态来。

\textbf{状态复用。} dp[i] 表示到达第 i 阶的方法数，转移是 dp[i]=dp[i-1]+dp[i-2]。再压缩成两个滚动变量。手算时只改被新输入影响的那一小部分，而不是回头重算完整候选。

\textbf{指针和字段。} 状态字段要能说清：prev2 表示 dp[i-2]，prev1 表示 dp[i-1]，current 表示 dp[i]。手算时按这个协议跟踪：n=4 时 dp[1]=1、dp[2]=2、dp[3]=3、dp[4]=5。每一步只依赖前两项。

\textbf{代码落点。} 关键是初始化：n=0 按题意通常返回 1 或不出现；LeetCode 输入 n>=1，n=1 返回 1，n=2 返回 2。关键代码行必须能逐句翻译成“旧状态怎样变成新状态”，否则就只是模板。

\textbf{正确性。} 所有到第 i 阶的方案按最后一步分成从 i-1 来和从 i-2 来两类，二者不重不漏。

\textbf{边界实验。} 先用n 为一、n 为二、题目是否允许零阶、结果范围做反例检查。实验时覆盖n 为一、n 为二、题目是否允许零阶、结果范围，先打印完整 DP 表，再比较空间压缩版本。若压缩版不同，优先检查遍历顺序和旧值保存。

\textbf{复杂度变化。} 暴力递归指数级；DP O(n) 时间，滚动变量 O(1) 空间。

\textbf{迁移追问。} 如果题目改成“若每次可走一到 k 阶，转移变成固定宽度窗口求和”，先判断原状态是否仍然覆盖新答案集合，再决定是微调边界、改 value 语义，还是换算法。

\subsection{LeetCode 198 打家劫舍}

\textbf{题目说明。} LeetCode 198 打家劫舍 的题面入口要先还原成具体任务：每间房只有偷和不偷两种结局，偷当前就不能偷前一间。

\textbf{样例入口。} 拿 [2, 7, 9] 手算，到房屋 9 时只有两类合法结局：不偷它，答案沿用前一间；偷它，就只能接前两间的最优。

\textbf{暴力基线。} 暴力递归每间房偷或不偷，偷了当前房就不能偷下一间，最后取最大金额。

\textbf{暴力过程。} 以 nums=[2, 7, 9] 为例，到房屋 9 时，合法结局只有两类：不偷 9 沿用前一间最优，偷 9 就只能接前两间最优。

\textbf{重复工作。} 题面模型：每间房只有偷和不偷两种结局，偷当前就不能偷前一间。暴力版的慢点要落到本题：浪费在递归中同一个下标之后的最优收益被反复计算。未来只取决于当前处理到哪一间。后面的优化只消掉这类重复工作，不能改变答案集合。

\textbf{优化条件。} 本题的关键观察是：滚动保存前一状态和前两状态。这句话必须能翻译成可维护状态；如果题目条件改变后观察不再成立，就不能继续套原来的写法。

\textbf{相邻状态。} 规律是 dp[i]=max(dp[i-1], dp[i-2]+nums[i])，相邻限制来自题面而不是公式本身。把这个特例继续拆开看：先问暴力搜索里哪些不同路径到达同一个后续问题，再给这个后续问题命名为状态。状态必须包含未来决策需要的全部信息，转移只从更小且已经算清的状态来。

\textbf{状态复用。} dp[i] 表示考虑前 i 间房的最大收益。转移为 max(dp[i-1], dp[i-2]+nums[i-1])，可压成两个变量。手算时只改被新输入影响的那一小部分，而不是回头重算完整候选。

\textbf{指针和字段。} 状态字段要能说清：previous\_one 是 dp[i-1]，previous\_two 是 dp[i-2]，money 是当前房屋金额。手算时按这个协议跟踪：[2, 7, 9] 中前两步最优到 7；处理 9 时比较不偷得 7、偷得 2+9=11，更新为 11。

\textbf{代码落点。} 关键是滚动更新顺序：先算 current=max(previous\_one, previous\_two+money)，再 previous\_two=previous\_one，previous\_one=current。关键代码行必须能逐句翻译成“旧状态怎样变成新状态”，否则就只是模板。

\textbf{正确性。} 最优方案对第 i 间房只有偷或不偷两类，二者互斥且覆盖全部合法方案。

\textbf{边界实验。} 先用空数组、单间、全零、金额很大做反例检查。实验时覆盖空数组、单间、全零、金额很大，先打印完整 DP 表，再比较空间压缩版本。若压缩版不同，优先检查遍历顺序和旧值保存。

\textbf{复杂度变化。} 暴力递归指数级；DP O(n) 时间，滚动 O(1) 空间。

\textbf{迁移追问。} 如果题目改成“环形房屋要拆成不含首和不含尾两个线性问题”，先判断原状态是否仍然覆盖新答案集合，再决定是微调边界、改 value 语义，还是换算法。

\subsection{LeetCode 322 零钱兑换}

\textbf{题目说明。} LeetCode 322 零钱兑换 的题面入口要先还原成具体任务：金额是容量，硬币可无限使用，目标是最少硬币数。

\textbf{样例入口。} 拿 coins=[1, 2, 5]、amount=11 手算，最后一步可以用 1、2 或 5，若用 5，则来源是 dp[6]+1。

\textbf{暴力基线。} 暴力递归枚举最后使用哪枚硬币，或者枚举每种硬币使用多少次，直到金额变成 0。

\textbf{暴力过程。} 以 coins=[1, 2, 5]、amount=11 为例，若最后一枚是 5，前面必须先凑出 6；若最后一枚是 2，前面必须先凑出 9。

\textbf{重复工作。} 题面模型：金额是容量，硬币可无限使用，目标是最少硬币数。暴力版的慢点要落到本题：浪费在同一个剩余金额会从很多硬币序列到达。凑出金额 x 的最少硬币数和之前路径无关，只和 x 有关。后面的优化只消掉这类重复工作，不能改变答案集合。

\textbf{优化条件。} 本题的关键观察是：dp[value] 从 value 减 coin 的已知状态转移。这句话必须能翻译成可维护状态；如果题目条件改变后观察不再成立，就不能继续套原来的写法。

\textbf{相邻状态。} 规律是 dp[x] 表示凑成 x 的最少硬币数，完全背包允许同一硬币重复使用；不可达状态不能参与加一。把这个特例继续拆开看：阶段是物品，容量是资源，状态值是可达、最值还是计数。一个物品能否在同一轮再次使用，直接决定容量循环方向；组合和排列也由循环顺序表达。

\textbf{状态复用。} dp[x] 表示凑出金额 x 的最少硬币数，初始化 dp[0]=0、其他为 INF。对每个金额尝试所有 coin。手算时只改被新输入影响的那一小部分，而不是回头重算完整候选。

\textbf{指针和字段。} 状态字段要能说清：x 是当前金额，coin 是最后一枚硬币，dp[x-coin]+1 是使用 coin 的候选答案。手算时按这个协议跟踪：dp[5] 可由 coin=5 得到 1；dp[11] 可由 dp[6]+1、dp[9]+1、dp[10]+1 取最小。

\textbf{代码落点。} 关键是不可达 INF 不能直接加一造成溢出；完全背包一维写法可金额正序，也可金额外层逐项求最小。关键代码行必须能逐句翻译成“旧状态怎样变成新状态”，否则就只是模板。

\textbf{正确性。} 任意最优方案都有一枚最后硬币 coin，去掉它后剩余金额 x-coin 必须也是最优子问题，否则可替换变更优。

\textbf{边界实验。} 先用无法凑出、amount 为零、硬币面额大于目标、哨兵溢出做反例检查。实验时覆盖无法凑出、amount 为零、硬币面额大于目标、哨兵溢出，并故意把容量方向写反构造最小反例。这样能确认循环方向来自题意，不是背模板。

\textbf{复杂度变化。} 暴力指数级；DP O(amount*coins.size()) 时间，O(amount) 空间。

\textbf{迁移追问。} 如果题目改成“零钱兑换 II 问组合数，循环语义不同”，先判断原状态是否仍然覆盖新答案集合，再决定是微调边界、改 value 语义，还是换算法。

\subsection{LeetCode 300 最长递增子序列}

\textbf{题目说明。} LeetCode 300 最长递增子序列 的题面入口要先还原成具体任务：二次 DP 固定结尾，二分优化维护每个长度的最小结尾。

\textbf{样例入口。} 拿 [10, 9, 2, 5, 3] 手算。二次 DP 固定每个位置为结尾，5 可接在 2 后；优化时 tails[长度] 保存该长度递增子序列的最小尾值，尾值越小未来越容易接。

\textbf{暴力基线。} 暴力递归或 DP 枚举每个元素作为子序列结尾，向前寻找所有更小元素。

\textbf{暴力过程。} 以 [10, 9, 2, 5, 3, 7, 101, 18] 为例，二次 DP 可以算出以 7 结尾的长度来自 2、5 或 3 后面接 7。

\textbf{重复工作。} 题面模型：二次 DP 固定结尾，二分优化维护每个长度的最小结尾。暴力版的慢点要落到本题：二次 DP 慢在每个位置都回看所有历史元素。更快写法不需要保存真实序列，只保存每个长度的最小可能结尾。后面的优化只消掉这类重复工作，不能改变答案集合。

\textbf{优化条件。} 本题的关键观察是：tails 不是实际答案序列，而是未来扩展潜力表。这句话必须能翻译成可维护状态；如果题目条件改变后观察不再成立，就不能继续套原来的写法。

\textbf{相邻状态。} 规律是二分替换 tails，不代表真实序列完整内容，而是维护未来潜力。把这个特例继续拆开看：先问暴力搜索里哪些不同路径到达同一个后续问题，再给这个后续问题命名为状态。状态必须包含未来决策需要的全部信息，转移只从更小且已经算清的状态来。

\textbf{状态复用。} tails[len] 表示长度为 len+1 的递增子序列中最小结尾。对每个 x，用 lower\_bound 找第一个 >=x 的位置并替换。手算时只改被新输入影响的那一小部分，而不是回头重算完整候选。

\textbf{指针和字段。} 状态字段要能说清：tails 不是最终答案序列，而是未来扩展潜力表；tails.size() 是当前最长长度。手算时按这个协议跟踪：处理 2、5、3 时，tails 从 [2, 5] 变成 [2, 3]，长度没变，但结尾更小，未来更容易接 7。

\textbf{代码落点。} 严格递增用 lower\_bound；若题目改成非递减，要改用 upper\_bound。要还原路径需要额外 predecessor 数组。关键代码行必须能逐句翻译成“旧状态怎样变成新状态”，否则就只是模板。

\textbf{正确性。} 同样长度下，结尾越小，未来能接上的数字集合只会更多；替换 tails 不改变已知长度，只提升扩展潜力。

\textbf{边界实验。} 先用重复元素、严格递增和非递减区别、空数组做反例检查。实验时覆盖重复元素、严格递增和非递减区别、空数组，先打印完整 DP 表，再比较空间压缩版本。若压缩版不同，优先检查遍历顺序和旧值保存。

\textbf{复杂度变化。} 二次 DP O(\ensuremath{n^{2}})，tails 加二分 O(n log n) 时间、O(n) 空间。

\textbf{迁移追问。} 如果题目改成“若要还原路径，需要额外前驱数组，不能只看 tails”，先判断原状态是否仍然覆盖新答案集合，再决定是微调边界、改 value 语义，还是换算法。

\subsection{LeetCode 1143 最长公共子序列}

\textbf{题目说明。} LeetCode 1143 最长公共子序列 的题面入口要先还原成具体任务：状态表示两个前缀的最长公共子序列长度。

\textbf{样例入口。} 拿 abcde 和 ace 手算，末尾 e 相等时答案来自两个前缀 abcd 和 ac 加一；若末尾不同，就取删掉一侧末尾后的较大值。

\textbf{暴力基线。} 暴力递归枚举两个字符串的子序列选择，或从末尾开始相等就选、不等就分支删一边。

\textbf{暴力过程。} 以 text1=abcde、text2=ace 为例，末尾 e 相等时，答案来自 abcd 和 ac 的 LCS 加一；若末尾不同，就只能丢掉一侧末尾试更优。

\textbf{重复工作。} 题面模型：状态表示两个前缀的最长公共子序列长度。暴力版的慢点要落到本题：浪费在很多前缀对会被重复计算。只要两个前缀长度 i、j 固定，它们的 LCS 长度就是固定状态。后面的优化只消掉这类重复工作，不能改变答案集合。

\textbf{优化条件。} 本题的关键观察是：末尾相同取左上加一，不同取上方和左方最大。这句话必须能翻译成可维护状态；如果题目条件改变后观察不再成立，就不能继续套原来的写法。

\textbf{相邻状态。} 规律是 dp[i][j] 表示两个前缀的 LCS 长度，转移只看左、上、左上。把这个特例继续拆开看：先问暴力搜索里哪些不同路径到达同一个后续问题，再给这个后续问题命名为状态。状态必须包含未来决策需要的全部信息，转移只从更小且已经算清的状态来。

\textbf{状态复用。} dp[i][j] 表示 text1 前 i 个字符和 text2 前 j 个字符的 LCS 长度。相等看左上加一，不等取上方和左方最大。手算时只改被新输入影响的那一小部分，而不是回头重算完整候选。

\textbf{指针和字段。} 状态字段要能说清：i/j 是前缀长度，dp[i-1][j] 表示丢 text1 末尾，dp[i][j-1] 表示丢 text2 末尾，dp[i-1][j-1] 表示两个末尾配对。手算时按这个协议跟踪：a 与 a 相等得到 1；到 abcde 和 ace 的最后，e 与 e 相等，所以由 abcd/ac 的答案加一。

\textbf{代码落点。} 关键是表多开一行一列表示空前缀，循环从 1 开始，访问字符用 text[i-1]。关键代码行必须能逐句翻译成“旧状态怎样变成新状态”，否则就只是模板。

\textbf{正确性。} 任意 LCS 对末尾字符只有配对或至少丢掉一边两类；转移覆盖这两类并取最大。

\textbf{边界实验。} 先用空串、完全相同、无公共字符、子序列不是子串做反例检查。实验时覆盖空串、完全相同、无公共字符、子序列不是子串，先打印完整 DP 表，再比较空间压缩版本。若压缩版不同，优先检查遍历顺序和旧值保存。

\textbf{复杂度变化。} 暴力指数级；DP O(mn) 时间和空间，可滚动到 O(min(m, n)) 空间。

\textbf{迁移追问。} 如果题目改成“若要求还原序列，需要从表右下角反向追踪选择”，先判断原状态是否仍然覆盖新答案集合，再决定是微调边界、改 value 语义，还是换算法。

\subsection{LeetCode 72 编辑距离}

\textbf{题目说明。} LeetCode 72 编辑距离 的题面入口要先还原成具体任务：二维状态表示两个前缀之间的最少编辑操作。

\textbf{样例入口。} 拿 word1=ab、word2=ac 手算最后一个字符。若 b 改成 c，是 dp[1][1]+1；若删掉 b，是 dp[1][2]+1；若插入 c，是 dp[2][1]+1。字符相等时才直接继承左上。

\textbf{暴力基线。} 暴力递归比较两个字符串后缀：相等就同时跳过，不等就尝试插入、删除、替换三种操作。

\textbf{暴力过程。} 以 word1=ab、word2=ac 为例，末尾 b 和 c 不同。替换 b 成 c 来自 dp[1][1]+1；删除 b 来自 dp[1][2]+1；插入 c 来自 dp[2][1]+1。

\textbf{重复工作。} 题面模型：二维状态表示两个前缀之间的最少编辑操作。暴力版的慢点要落到本题：浪费在不同编辑序列会反复到达同一对前缀长度。只要 i 和 j 相同，后续最少编辑次数就相同。后面的优化只消掉这类重复工作，不能改变答案集合。

\textbf{优化条件。} 本题的关键观察是：末尾字符相同继承左上；不同则枚举插入、删除、替换。这句话必须能翻译成可维护状态；如果题目条件改变后观察不再成立，就不能继续套原来的写法。

\textbf{相邻状态。} 规律是每个格子枚举最后一次编辑操作，而不是凭感觉填三邻居最小。把这个特例继续拆开看：先问暴力搜索里哪些不同路径到达同一个后续问题，再给这个后续问题命名为状态。状态必须包含未来决策需要的全部信息，转移只从更小且已经算清的状态来。

\textbf{状态复用。} dp[i][j] 表示 word1 前 i 个字符变成 word2 前 j 个字符的最少操作数。按前缀长度从小到大填表。手算时只改被新输入影响的那一小部分，而不是回头重算完整候选。

\textbf{指针和字段。} 状态字段要能说清：i/j 是前缀长度；dp[i-1][j] 对应删除，dp[i][j-1] 对应插入，dp[i-1][j-1] 对应替换或字符相等继承。手算时按这个协议跟踪：空串到 ac 需要插入 2 次；ab 到空串需要删除 2 次。内部格子按末尾字符是否相等决定转移。

\textbf{代码落点。} 关键是第一行 dp[0][j]=j、第一列 dp[i][0]=i。字符相等时不要额外加一。关键代码行必须能逐句翻译成“旧状态怎样变成新状态”，否则就只是模板。

\textbf{正确性。} 任意最优编辑序列的最后一步只有插入、删除、替换或不操作四类，转移枚举了全部最后一步。

\textbf{边界实验。} 先用空串、完全相同、完全不同、操作代价不同做反例检查。实验时覆盖空串、完全相同、完全不同、操作代价不同，先打印完整 DP 表，再比较空间压缩版本。若压缩版不同，优先检查遍历顺序和旧值保存。

\textbf{复杂度变化。} 暴力递归指数级；二维 DP O(mn) 时间和空间，可滚动到 O(min(m, n)) 空间。

\textbf{迁移追问。} 如果题目改成“若只允许插入删除，问题退化到和最长公共子序列相关”，先判断原状态是否仍然覆盖新答案集合，再决定是微调边界、改 value 语义，还是换算法。

\subsection{LeetCode 416 分割等和子集}

\textbf{题目说明。} LeetCode 416 分割等和子集 的题面入口要先还原成具体任务：总和一半是背包容量，每个数只能使用一次。

\textbf{样例入口。} 拿 [1, 5, 11, 5] 手算，总和 22，目标容量 11；可以用 11 单独凑成，也可以 5+5+1。

\textbf{暴力基线。} 暴力枚举每个数放到左集合还是右集合，最后检查两边和是否相等。

\textbf{暴力过程。} 以 [1, 5, 11, 5] 为例，总和 22，目标是找一个子集和为 11；可以选单独的 11，也可以选 5+5+1。

\textbf{重复工作。} 题面模型：总和一半是背包容量，每个数只能使用一次。暴力版的慢点要落到本题：浪费在不同选择顺序会反复得到相同剩余容量。问题只关心能否凑到 sum/2，不关心具体左右集合顺序。后面的优化只消掉这类重复工作，不能改变答案集合。

\textbf{优化条件。} 本题的关键观察是：0-1 背包可达性，一维容量必须倒序。这句话必须能翻译成可维护状态；如果题目条件改变后观察不再成立，就不能继续套原来的写法。

\textbf{相邻状态。} 规律是分割等和转成 0-1 背包是否能凑到 sum/2，容量必须倒序避免重复使用同一元素。把这个特例继续拆开看：阶段是物品，容量是资源，状态值是可达、最值还是计数。一个物品能否在同一轮再次使用，直接决定容量循环方向；组合和排列也由循环顺序表达。

\textbf{状态复用。} 先判断总和是否为偶数。dp[cap] 表示处理过若干数字后是否能凑出 cap，每个 num 只能用一次，所以容量倒序更新。手算时只改被新输入影响的那一小部分，而不是回头重算完整候选。

\textbf{指针和字段。} 状态字段要能说清：target=sum/2，cap 是容量，num 是当前物品；dp[cap-num] 为真则 dp[cap] 可变真。手算时按这个协议跟踪：处理 1 后可达 1；处理 5 后可达 5 和 6；处理 11 后 target 11 直接可达。

\textbf{代码落点。} 关键是 for cap from target down to num。若正序更新，会在同一轮把当前 num 用多次，变成完全背包。关键代码行必须能逐句翻译成“旧状态怎样变成新状态”，否则就只是模板。

\textbf{正确性。} 每个数只有选或不选两种状态；倒序容量保证转移读取的是上一阶段结果，不会重复使用当前数。

\textbf{边界实验。} 先用总和为奇数、目标为零、数值较大、复杂度依赖 target做反例检查。实验时覆盖总和为奇数、目标为零、数值较大、复杂度依赖 target，并故意把容量方向写反构造最小反例。这样能确认循环方向来自题意，不是背模板。

\textbf{复杂度变化。} 暴力 O(\ensuremath{2^{n}})，0-1 背包 O(n*target) 时间、O(target) 空间。

\textbf{迁移追问。} 如果题目改成“负数会破坏普通容量模型，需要重新建模”，先判断原状态是否仍然覆盖新答案集合，再决定是微调边界、改 value 语义，还是换算法。

\subsection{LeetCode 494 目标和}

\textbf{题目说明。} LeetCode 494 目标和 的题面入口要先还原成具体任务：正负号选择可代数转换为选若干数凑出特定正和。

\textbf{样例入口。} 拿 [1, 1, 1, 1, 1]、target=3 手算，给其中一个 1 加负号，其余为正可以得到 3，共有 5 种位置选择。也可转成选正数子集和为 (sum+target)/2。

\textbf{暴力基线。} 暴力递归给每个数添加 + 或 -，枚举 \ensuremath{2^{n}} 个符号序列并统计和为 target 的方案。

\textbf{暴力过程。} 以 [1, 1, 1, 1, 1]、target=3 为例，只有一个 1 取负、其余取正时可得到 3，共 5 种位置选择。

\textbf{重复工作。} 题面模型：正负号选择可代数转换为选若干数凑出特定正和。暴力版的慢点要落到本题：浪费在符号序列会反复得到相同部分和。把正数集合 P、负数集合 N 写成 sum(P)-sum(N)=target，可转成子集和计数。后面的优化只消掉这类重复工作，不能改变答案集合。

\textbf{优化条件。} 本题的关键观察是：若 S 加 target 为奇数或目标绝对值过大，直接无解。这句话必须能翻译成可维护状态；如果题目条件改变后观察不再成立，就不能继续套原来的写法。

\textbf{相邻状态。} 规律是符号分配变成 0-1 计数背包，奇偶和范围先判定。把这个特例继续拆开看：阶段是物品，容量是资源，状态值是可达、最值还是计数。一个物品能否在同一轮再次使用，直接决定容量循环方向；组合和排列也由循环顺序表达。

\textbf{状态复用。} 总和 total。若 total+target 为奇数或 abs(target)>total，直接 0；否则找子集和 positive=(total+target)/2 的方案数。手算时只改被新输入影响的那一小部分，而不是回头重算完整候选。

\textbf{指针和字段。} 状态字段要能说清：dp[cap] 是凑出 cap 的方案数，num 是当前数；0-1 计数背包容量倒序更新。手算时按这个协议跟踪：target=3、total=5，positive=4；选择四个 1 作为正数等价于选择哪个 1 不放进去，共 5 种。

\textbf{代码落点。} 关键是 dp[0]=1，for cap from positive down to num。数字 0 会让方案数翻倍，倒序更新能自然处理。关键代码行必须能逐句翻译成“旧状态怎样变成新状态”，否则就只是模板。

\textbf{正确性。} 每个符号分配和一个正数子集一一对应；子集和满足 positive 时，正负差正好为 target。

\textbf{边界实验。} 先用零元素会让方案数翻倍、目标为负、容量为零做反例检查。实验时覆盖零元素会让方案数翻倍、目标为负、容量为零，并故意把容量方向写反构造最小反例。这样能确认循环方向来自题意，不是背模板。

\textbf{复杂度变化。} 暴力 O(\ensuremath{2^{n}})，背包 O(n*positive) 时间，O(positive) 空间。

\textbf{迁移追问。} 如果题目改成“也可用哈希 DP 保存所有可能和，但背包转换更高效”，先判断原状态是否仍然覆盖新答案集合，再决定是微调边界、改 value 语义，还是换算法。

\subsection{LeetCode 55 跳跃游戏}

\textbf{题目说明。} LeetCode 55 跳跃游戏 的题面入口要先还原成具体任务：扫描过程中只关心当前能到达的最远位置。

\textbf{样例入口。} 拿 [2, 3, 1, 1, 4] 手算，下标 0 能覆盖到 2；走到下标 1 时最远可达更新到 4，已经覆盖终点。再拿 [0, 1]，扫描到下标 1 时已经超过最远可达，失败。

\textbf{暴力基线。} 暴力把每个下标能跳到的位置都当成分支搜索，尝试是否存在一条路径能到终点。

\textbf{暴力过程。} 以 [2, 3, 1, 1, 4] 为例，从 0 可以跳到 1 或 2；若到 1，最远能继续覆盖到 4。以 [0, 1] 为例，从 0 一步也走不出去。

\textbf{重复工作。} 题面模型：扫描过程中只关心当前能到达的最远位置。暴力版的慢点要落到本题：浪费在枚举具体路径。判断能否到达终点时，不关心走的是哪条路径，只关心目前所有可达位置能覆盖到的最远下标。后面的优化只消掉这类重复工作，不能改变答案集合。

\textbf{优化条件。} 本题的关键观察是：如果下标超过最远可达，任何之前路径都不能到达它；否则用它扩展边界。这句话必须能翻译成可维护状态；如果题目条件改变后观察不再成立，就不能继续套原来的写法。

\textbf{相邻状态。} 规律是只维护当前能到达的最远边界，任何超过边界的位置都不可能由前面路径到达。把这个特例继续拆开看：先构造一个非贪心选择，再比较两者留下的资源、右边界或剩余时间。只有能把非贪心第一步交换成贪心第一步且不变差，局部选择才是规律。

\textbf{状态复用。} 扫描下标 i，维护 farthest。若 i>farthest，说明当前下标不可达，直接失败；否则用 i+nums[i] 更新 farthest。手算时只改被新输入影响的那一小部分，而不是回头重算完整候选。

\textbf{指针和字段。} 状态字段要能说清：i 是当前检查位置，farthest 是所有已可达位置能跳到的最远边界，n-1 是目标边界。手算时按这个协议跟踪：样例中 i=0 后 farthest=2；i=1 可达并把 farthest 更新为 4，已经覆盖终点。

\textbf{代码落点。} 关键顺序是先判断当前位置是否超过 farthest；若没超过，再用 i+nums[i] 扩展最远可达位置。数组长度为 1 时起点就是终点。关键代码行必须能逐句翻译成“旧状态怎样变成新状态”，否则就只是模板。

\textbf{正确性。} 在扫描到 i 前，farthest 汇总了所有可达下标的最远覆盖；若 i 超过它，任何旧路径都到不了 i，更不可能通过 i 之后的位置到终点。

\textbf{边界实验。} 先用第一个元素为零、数组长度一、恰好到达最后、远超最后做反例检查。实验时除了第一个元素为零、数组长度一、恰好到达最后、远超最后，还要故意替换成一个看似合理但错误的局部规则，构造反例。能解释错误贪心为何失败，才算理解正确贪心。

\textbf{复杂度变化。} 路径搜索最坏指数级；贪心边界扫描 O(n) 时间、O(1) 空间。

\textbf{迁移追问。} 如果题目改成“求最少跳数时把可达区间看成 BFS 层”，先判断原状态是否仍然覆盖新答案集合，再决定是微调边界、改 value 语义，还是换算法。

\subsection{LeetCode 134 加油站}

\textbf{题目说明。} LeetCode 134 加油站 的题面入口要先还原成具体任务：绕行一圈是否可行取决于总油量差和局部剩余油量。

\textbf{样例入口。} 拿 gas=[1, 2, 3, 4, 5]、cost=[3, 4, 5, 1, 2] 手算，从 0 出发中途油量为负，0 到失败点之间任何站都不能作为起点；从 3 出发可以绕完。

\textbf{暴力基线。} 暴力把每个加油站都当起点，模拟绕一圈，若油量中途为负就换下一个起点。

\textbf{暴力过程。} gas=[1, 2, 3, 4, 5]、cost=[3, 4, 5, 1, 2] 中，从 0 出发会在中途油量为负；从 3 出发可以完成一圈。

\textbf{重复工作。} 题面模型：绕行一圈是否可行取决于总油量差和局部剩余油量。暴力版的慢点要落到本题：浪费在失败起点之后仍逐个重试中间站。若从 start 到 i 的累计油量已经为负，那么 start..i 中任何站都不可能作为可行起点。后面的优化只消掉这类重复工作，不能改变答案集合。

\textbf{优化条件。} 本题的关键观察是：总和为负必不可行；扫描时若当前剩余为负，起点只能放到下一站。这句话必须能翻译成可维护状态；如果题目条件改变后观察不再成立，就不能继续套原来的写法。

\textbf{相邻状态。} 规律是总油量差为负则无解，局部剩余为负时起点跳到下一站。把这个特例继续拆开看：先构造一个非贪心选择，再比较两者留下的资源、右边界或剩余时间。只有能把非贪心第一步交换成贪心第一步且不变差，局部选择才是规律。

\textbf{状态复用。} 扫描一次，total 记录总油量差，tank 记录当前候选起点到当前位置的剩余油量，start 是当前候选起点。手算时只改被新输入影响的那一小部分，而不是回头重算完整候选。

\textbf{指针和字段。} 状态字段要能说清：diff=gas[i]-cost[i]，tank 是局部剩余，total 是全局可行性，start 是下一次尝试起点。手算时按这个协议跟踪：样例前几站累计 tank 变负后，把 start 移到下一站并清零 tank；最后 total>=0，返回最终 start=3。

\textbf{代码落点。} 关键是 total 全程累加，tank<0 时 start=i+1、tank=0。循环结束后 total<0 返回 -1，否则返回 start。关键代码行必须能逐句翻译成“旧状态怎样变成新状态”，否则就只是模板。

\textbf{正确性。} tank<0 说明从当前 start 到 i 的总补给不足；其中任意中间站作为起点，少了前面一段补给，只会更差，故可整体跳过。

\textbf{边界实验。} 先用总和刚好为零、多个可行起点、某站油量为零、消耗大于补给做反例检查。实验时除了总和刚好为零、多个可行起点、某站油量为零、消耗大于补给，还要故意替换成一个看似合理但错误的局部规则，构造反例。能解释错误贪心为何失败，才算理解正确贪心。

\textbf{复杂度变化。} 暴力 O(\ensuremath{n^{2}})，一次扫描 O(n) 时间、O(1) 空间。

\textbf{迁移追问。} 如果题目改成“它的领先性证明是前一段任意站点都不能作为起点”，先判断原状态是否仍然覆盖新答案集合，再决定是微调边界、改 value 语义，还是换算法。

\subsection{LeetCode 56 合并区间}

\textbf{题目说明。} LeetCode 56 合并区间 的题面入口要先还原成具体任务：排序后可能与当前区间相交的只会是结果尾部。

\textbf{样例入口。} 拿 [[1, 3], [2, 6], [8, 10]] 手算，排序后 [2, 6] 只可能和结果尾部 [1, 3] 相交，合并成 [1, 6]；[8, 10] 的起点大于尾部右端，另开新区间。

\textbf{暴力基线。} 暴力可以不断挑一个区间，和所有其他区间比较是否相交，能合并就合并，直到没有变化。

\textbf{暴力过程。} 以 [[1, 3], [2, 6], [8, 10]] 为例，[1, 3] 和 [2, 6] 相交，应合成 [1, 6]；[8, 10] 与它不相交，单独保留。

\textbf{重复工作。} 题面模型：排序后可能与当前区间相交的只会是结果尾部。暴力版的慢点要落到本题：浪费在反复两两比较。按左端排序后，能和当前结果尾部相交的区间会连续出现，不需要再回头看更早区间。后面的优化只消掉这类重复工作，不能改变答案集合。

\textbf{优化条件。} 本题的关键观察是：按起点排序，扫描时维护结果最后区间的右端点。这句话必须能翻译成可维护状态；如果题目条件改变后观察不再成立，就不能继续套原来的写法。

\textbf{相邻状态。} 规律是排序后当前区间只需要看结果尾部，端点相接是否合并要按闭区间语义决定。把这个特例继续拆开看：排序以后，真正还能影响当前区间的候选必须贴近当前端点、结果尾部或资源堆顶。某个区间一旦被覆盖、合并或弹出，就要说明它为什么不会和后续区间重新产生更优关系。

\textbf{状态复用。} 先按 start 升序排序。结果数组保持已合并且不重叠；当前区间只和结果最后一个区间比较。手算时只改被新输入影响的那一小部分，而不是回头重算完整候选。

\textbf{指针和字段。} 状态字段要能说清：last 是结果尾区间，interval 是当前扫描区间；若 interval.start<=last.end 就合并，否则追加新区间。手算时按这个协议跟踪：排序后先放 [1, 3]；看到 [2, 6]，2<=3，更新尾部右端为 6；看到 [8, 10]，8>6，追加。

\textbf{代码落点。} 关键是明确闭区间语义，端点相接是否合并由题目决定。本题闭区间重叠时用 start<=last.end。关键代码行必须能逐句翻译成“旧状态怎样变成新状态”，否则就只是模板。

\textbf{正确性。} 排序保证当前区间之前的所有区间左端不大于它；若它不和结果尾相交，就不可能和更早且右端不超过尾部的区间相交。

\textbf{边界实验。} 先用端点相接、完全覆盖、无重叠、空列表、输入未排序做反例检查。实验时覆盖端点相接、完全覆盖、无重叠、空列表、输入未排序，画出排序后的区间序列和当前结果尾部。动态版本要专门测前驱、后继和端点相接。

\textbf{复杂度变化。} 暴力反复合并最坏 O(\ensuremath{n^{2}})，排序扫描 O(n log n) 时间，输出空间 O(n)。

\textbf{迁移追问。} 如果题目改成“若区间是半开区间，端点相等不一定合并”，先判断原状态是否仍然覆盖新答案集合，再决定是微调边界、改 value 语义，还是换算法。

\subsection{LeetCode 684 冗余连接}

\textbf{题目说明。} LeetCode 684 冗余连接 的题面入口要先还原成具体任务：树加一条边会形成唯一环。

\textbf{样例入口。} 拿 edges=[[1, 2], [1, 3], [2, 3]] 手算，前两条边把 1、2、3 连成一棵树，第三条边连接的两个点已经同组，所以它就是冗余边。

\textbf{暴力基线。} 慢版本每次合并后用 DFS 或 BFS 重新判断连通性。它正确但重复遍历已有连接；当关系只会合并不会拆开时，可以压成集合代表。放到本题就是：先按“树加一条边会形成唯一环”完整试慢版本；样例里已经能看到“规律是无向图加边时，第一次 union 失败的边形成环”，所以优化前必须先能说清慢版本枚举了哪些候选。

\textbf{暴力过程。} 拿 edges=[[1, 2], [1, 3], [2, 3]] 手算，前两条边把 1、2、3 连成一棵树，第三条边连接的两个点已经同组，所以它就是冗余边。

\textbf{重复工作。} 题面模型：树加一条边会形成唯一环。暴力版的慢点要落到本题：慢版本会围绕“树加一条边会形成唯一环”反复搜索连通性。样例规律是“规律是无向图加边时，第一次 union 失败的边形成环”，说明关系只需要被压成集合代表；每次 union 失败或成功都要能翻译回题意。后面的优化只消掉这类重复工作，不能改变答案集合。

\textbf{优化条件。} 本题的关键观察是：按顺序加边，若一条边两端已连通，则它就是冗余边。这句话必须能翻译成可维护状态；如果题目条件改变后观察不再成立，就不能继续套原来的写法。

\textbf{相邻状态。} 规律是无向图加边时，第一次 union 失败的边形成环。把这个特例继续拆开看：union 只是在维护等价类，不是在保存具体路径。两个节点代表元相同只能说明连通或关系一致；如果题目还要距离、比例或奇偶性，必须把附加信息挂到根关系上。

\textbf{状态复用。} 把连通分量压成代表元；查询只比较代表，合并只发生在两个根之间。本题落点是：规律是无向图加边时，第一次 union 失败的边形成环。手算时只改被新输入影响的那一小部分，而不是回头重算完整候选。

\textbf{指针和字段。} 状态字段要能说清：parent、size 或 rank、find 返回的代表元、合并时的附加信息；带权或奇偶并查集还要维护到根的关系。手算时按这个协议跟踪：拿 edges=[[1, 2], [1, 3], [2, 3]] 手算，前两条边把 1、2、3 连成一棵树，第三条边连接的两个点已经同组，所以它就是冗余边。然后按协议复查：手算时写 parent 表和集合代表。每次 union 只改变根之间的关系，find 后代表相同才说明连通。

\textbf{代码落点。} find 路径压缩不能破坏附加权值；union 只能连接两个根；合并失败和重复边要保留题意解释。关键代码行必须能逐句翻译成“旧状态怎样变成新状态”，否则就只是模板。

\textbf{正确性。} 证明从等价关系开始：合并维护连通性的传递性，find 返回的代表元就是当前集合身份。

\textbf{边界实验。} 先用节点编号、最后一条成环边、重复边、并查集初始化大小做反例检查。实验时覆盖节点编号、最后一条成环边、重复边、并查集初始化大小，打印每个元素的代表元和集合大小。重复合并、已经连通的边和字符串编号最容易暴露实现问题。

\textbf{复杂度变化。} 暴力每次查询都搜索连通性代价高；并查集把合并和查询摊还到近似常数，整体接近线性。

\textbf{迁移追问。} 如果题目改成“有向冗余连接需要同时处理入度为二和环，复杂度更高”，先判断原状态是否仍然覆盖新答案集合，再决定是微调边界、改 value 语义，还是换算法。

\subsection{LeetCode 721 账户合并}

\textbf{题目说明。} LeetCode 721 账户合并 的题面入口要先还原成具体任务：共享邮箱说明账户属于同一人，邮箱之间形成连通分量。

\textbf{样例入口。} 拿 John 的账号 [a, b] 和另一个 [b, c] 手算，邮箱 b 重叠，所以 a、b、c 属于同一人；名字只是输出标签，合并依据是邮箱。

\textbf{暴力基线。} 慢版本每次合并后用 DFS 或 BFS 重新判断连通性。它正确但重复遍历已有连接；当关系只会合并不会拆开时，可以压成集合代表。放到本题就是：先按“共享邮箱说明账户属于同一人，邮箱之间形成连通分量”完整试慢版本；样例里已经能看到“规律是邮箱映射到编号后并查集合并，最后按代表收集并排序邮箱”，所以优化前必须先能说清慢版本枚举了哪些候选。

\textbf{暴力过程。} 拿 John 的账号 [a, b] 和另一个 [b, c] 手算，邮箱 b 重叠，所以 a、b、c 属于同一人；名字只是输出标签，合并依据是邮箱。

\textbf{重复工作。} 题面模型：共享邮箱说明账户属于同一人，邮箱之间形成连通分量。暴力版的慢点要落到本题：慢版本会围绕“共享邮箱说明账户属于同一人，邮箱之间形成连通分量”反复搜索连通性。样例规律是“规律是邮箱映射到编号后并查集合并，最后按代表收集并排序邮箱”，说明关系只需要被压成集合代表；每次 union 失败或成功都要能翻译回题意。后面的优化只消掉这类重复工作，不能改变答案集合。

\textbf{优化条件。} 本题的关键观察是：给邮箱编号并查集合并，同根邮箱收集后排序输出。这句话必须能翻译成可维护状态；如果题目条件改变后观察不再成立，就不能继续套原来的写法。

\textbf{相邻状态。} 规律是邮箱映射到编号后并查集合并，最后按代表收集并排序邮箱。把这个特例继续拆开看：union 只是在维护等价类，不是在保存具体路径。两个节点代表元相同只能说明连通或关系一致；如果题目还要距离、比例或奇偶性，必须把附加信息挂到根关系上。

\textbf{状态复用。} 把连通分量压成代表元；查询只比较代表，合并只发生在两个根之间。本题落点是：规律是邮箱映射到编号后并查集合并，最后按代表收集并排序邮箱。手算时只改被新输入影响的那一小部分，而不是回头重算完整候选。

\textbf{指针和字段。} 状态字段要能说清：parent、size 或 rank、find 返回的代表元、合并时的附加信息；带权或奇偶并查集还要维护到根的关系。手算时按这个协议跟踪：拿 John 的账号 [a, b] 和另一个 [b, c] 手算，邮箱 b 重叠，所以 a、b、c 属于同一人；名字只是输出标签，合并依据是邮箱。然后按协议复查：手算时写 parent 表和集合代表。每次 union 只改变根之间的关系，find 后代表相同才说明连通。

\textbf{代码落点。} find 路径压缩不能破坏附加权值；union 只能连接两个根；合并失败和重复边要保留题意解释。关键代码行必须能逐句翻译成“旧状态怎样变成新状态”，否则就只是模板。

\textbf{正确性。} 证明从等价关系开始：合并维护连通性的传递性，find 返回的代表元就是当前集合身份。

\textbf{边界实验。} 先用同名不同人、一个账户多个邮箱、重复邮箱、输出排序做反例检查。实验时覆盖同名不同人、一个账户多个邮箱、重复邮箱、输出排序，打印每个元素的代表元和集合大小。重复合并、已经连通的边和字符串编号最容易暴露实现问题。

\textbf{复杂度变化。} 暴力每次查询都搜索连通性代价高；并查集把合并和查询摊还到近似常数，整体接近线性。

\textbf{迁移追问。} 如果题目改成“姓名不能作为合并依据，只能作为输出字段”，先判断原状态是否仍然覆盖新答案集合，再决定是微调边界、改 value 语义，还是换算法。

\subsection{LeetCode 208 实现 Trie}

\textbf{题目说明。} LeetCode 208 实现 Trie 的题面入口要先还原成具体任务：前缀树把字符串公共前缀共享成路径。

\textbf{样例入口。} 拿插入 apple 再查 app 手算，app 的路径存在，但如果 p 节点没有终止标记，search(app) 应为假，startsWith(app) 才为真。

\textbf{暴力基线。} 慢版本在每个节点重新扫描子树或路径，先求出局部答案。重复扫描暴露了真正状态：每个节点只需要从孩子或父亲那里拿到少量信息。放到本题就是：先按“前缀树把字符串公共前缀共享成路径”完整试慢版本；样例里已经能看到“规律是 Trie 节点既要有孩子边，也要有单词终止标记”，所以优化前必须先能说清慢版本枚举了哪些候选。

\textbf{暴力过程。} 拿插入 apple 再查 app 手算，app 的路径存在，但如果 p 节点没有终止标记，search(app) 应为假，startsWith(app) 才为真。

\textbf{重复工作。} 题面模型：前缀树把字符串公共前缀共享成路径。暴力版的慢点要落到本题：慢版本会在树上重复确认“前缀树把字符串公共前缀共享成路径”。样例规律是“规律是 Trie 节点既要有孩子边，也要有单词终止标记”，说明父子之间真正需要传递的是高度、上下界、命中状态、层号或两类选择值，而不是反复扫描整棵子树。后面的优化只消掉这类重复工作，不能改变答案集合。

\textbf{优化条件。} 本题的关键观察是：插入和查询都逐字符沿边移动，不存在的边按需要创建。这句话必须能翻译成可维护状态；如果题目条件改变后观察不再成立，就不能继续套原来的写法。

\textbf{相邻状态。} 规律是 Trie 节点既要有孩子边，也要有单词终止标记。把这个特例继续拆开看：节点向父亲返回的量和全局答案通常不是同一个量。先写叶子或空节点的返回值，再写父节点如何合并左右孩子，递归式才有边界基础。

\textbf{状态复用。} 把对子树的重复扫描压成返回值或队列状态；每个节点只合并孩子给出的稳定信息。本题落点是：规律是 Trie 节点既要有孩子边，也要有单词终止标记。手算时只改被新输入影响的那一小部分，而不是回头重算完整候选。

\textbf{指针和字段。} 状态字段要能说清：节点指针、传入约束或返回值、当前答案、层号或路径状态；空节点的返回值必须和状态定义匹配。手算时按这个协议跟踪：拿插入 apple 再查 app 手算，app 的路径存在，但如果 p 节点没有终止标记，search(app) 应为假，startsWith(app) 才为真。然后按协议复查：手算时从叶子开始写返回值，或者从根开始写传入约束。不要只写访问顺序，要写每条边上传递的信息。

\textbf{代码落点。} 递归版本先处理空节点；迭代版本的栈元素要带完整状态；全负值、重复值和极端深树要单独测。关键代码行必须能逐句翻译成“旧状态怎样变成新状态”，否则就只是模板。

\textbf{正确性。} 证明从子树归纳或层序不变量开始：孩子状态正确时父节点合并正确；若用队列，当前轮队列必须正好保存当前层节点。

\textbf{边界实验。} 先用空字符串约定、单词结束标记、前缀存在不等于单词存在做反例检查。实验时覆盖空字符串约定、单词结束标记、前缀存在不等于单词存在，尤其关注空节点、单节点、链式树和全负值。递归版本和显式栈版本可以互相对拍。

\textbf{复杂度变化。} 暴力在每个节点重复扫描子树或反复寻找层边界可能退化到平方级；后序汇总、显式栈或层序遍历让每个节点处理常数次，通常是线性时间。

\textbf{迁移追问。} 如果题目改成“若字符集很大，用哈希子节点；小写字母用数组更快”，先判断原状态是否仍然覆盖新答案集合，再决定是微调边界、改 value 语义，还是换算法。

\subsection{LeetCode 1319 连通网络的操作次数}

\textbf{题目说明。} LeetCode 1319 连通网络的操作次数 的题面入口要先还原成具体任务：多余网线可以挪用来连接不同连通分量。

\textbf{样例入口。} 拿 n=4、连接 [[0, 1], [0, 2], [1, 2]] 手算，前三个点已有一条冗余边，但节点 3 孤立；总边数只有 3，刚好可以拿冗余边连上 3。

\textbf{暴力基线。} 慢版本每次合并后用 DFS 或 BFS 重新判断连通性。它正确但重复遍历已有连接；当关系只会合并不会拆开时，可以压成集合代表。放到本题就是：先按“多余网线可以挪用来连接不同连通分量”完整试慢版本；样例里已经能看到“规律是边数少于 n-1 直接失败，否则答案是连通分量数减一”，所以优化前必须先能说清慢版本枚举了哪些候选。

\textbf{暴力过程。} 拿 n=4、连接 [[0, 1], [0, 2], [1, 2]] 手算，前三个点已有一条冗余边，但节点 3 孤立；总边数只有 3，刚好可以拿冗余边连上 3。

\textbf{重复工作。} 题面模型：多余网线可以挪用来连接不同连通分量。暴力版的慢点要落到本题：慢版本会围绕“多余网线可以挪用来连接不同连通分量”反复搜索连通性。样例规律是“规律是边数少于 n-1 直接失败，否则答案是连通分量数减一”，说明关系只需要被压成集合代表；每次 union 失败或成功都要能翻译回题意。后面的优化只消掉这类重复工作，不能改变答案集合。

\textbf{优化条件。} 本题的关键观察是：先判断边数至少为 n 减一，再用并查集统计连通分量，答案是分量数减一。这句话必须能翻译成可维护状态；如果题目条件改变后观察不再成立，就不能继续套原来的写法。

\textbf{相邻状态。} 规律是边数少于 n-1 直接失败，否则答案是连通分量数减一。把这个特例继续拆开看：union 只是在维护等价类，不是在保存具体路径。两个节点代表元相同只能说明连通或关系一致；如果题目还要距离、比例或奇偶性，必须把附加信息挂到根关系上。

\textbf{状态复用。} 把连通分量压成代表元；查询只比较代表，合并只发生在两个根之间。本题落点是：规律是边数少于 n-1 直接失败，否则答案是连通分量数减一。手算时只改被新输入影响的那一小部分，而不是回头重算完整候选。

\textbf{指针和字段。} 状态字段要能说清：parent、size 或 rank、find 返回的代表元、合并时的附加信息；带权或奇偶并查集还要维护到根的关系。手算时按这个协议跟踪：拿 n=4、连接 [[0, 1], [0, 2], [1, 2]] 手算，前三个点已有一条冗余边，但节点 3 孤立；总边数只有 3，刚好可以拿冗余边连上 3。然后按协议复查：手算时写 parent 表和集合代表。每次 union 只改变根之间的关系，find 后代表相同才说明连通。

\textbf{代码落点。} find 路径压缩不能破坏附加权值；union 只能连接两个根；合并失败和重复边要保留题意解释。关键代码行必须能逐句翻译成“旧状态怎样变成新状态”，否则就只是模板。

\textbf{正确性。} 证明从等价关系开始：合并维护连通性的传递性，find 返回的代表元就是当前集合身份。

\textbf{边界实验。} 先用边数不足、已经连通、重复边、多分量做反例检查。实验时覆盖边数不足、已经连通、重复边、多分量，打印每个元素的代表元和集合大小。重复合并、已经连通的边和字符串编号最容易暴露实现问题。

\textbf{复杂度变化。} 暴力每次查询都搜索连通性代价高；并查集把合并和查询摊还到近似常数，整体接近线性。

\textbf{迁移追问。} 如果题目改成“合并失败的边代表冗余资源，但只要总边数足够即可”，先判断原状态是否仍然覆盖新答案集合，再决定是微调边界、改 value 语义，还是换算法。

\subsection{LeetCode 127 单词接龙}

\textbf{题目说明。} LeetCode 127 单词接龙 的题面入口要先还原成具体任务：每个单词是节点，一次修改一个字符形成边。

\textbf{样例入口。} 拿 hit -> hot -> dot -> dog -> cog 手算，每次只能改一个字符，所以 BFS 第一层是 hot，第二层才到 dot 或 lot。

\textbf{暴力基线。} 暴力可以把每个单词和所有其他单词比较一遍，差一个字符就连边，再从 beginWord 做 BFS。

\textbf{暴力过程。} 以 hit 到 cog 为例，hit 只能一步到 hot；hot 下一层到 dot 或 lot。第一次到 cog 的层数就是最短转换长度。

\textbf{重复工作。} 题面模型：每个单词是节点，一次修改一个字符形成边。暴力版的慢点要落到本题：浪费在找邻居时每次扫描整个词表。单词长度固定时，可以把 h*t、*ot 这类通配模式作为桶，快速找到只差一位的邻居。后面的优化只消掉这类重复工作，不能改变答案集合。

\textbf{优化条件。} 本题的关键观察是：通配模式桶加速邻居查询，BFS 第一次到达终点就是最短转换长度。这句话必须能翻译成可维护状态；如果题目条件改变后观察不再成立，就不能继续套原来的写法。

\textbf{相邻状态。} 规律是单词作为节点，通配桶如 h*t 把邻居快速找出；第一次到达终点就是最短转换长度。把这个特例继续拆开看：状态节点不一定等于原始位置，还可能包含钥匙、剩余资源、时间或访问集合。visited 只能合并未来能力完全相同的状态，否则会把合法路径剪掉。

\textbf{状态复用。} 预处理 pattern -> words 的映射，BFS 队列保存当前单词和步数，visited 防止重复入队。手算时只改被新输入影响的那一小部分，而不是回头重算完整候选。

\textbf{指针和字段。} 状态字段要能说清：word 是当前单词，pattern 是替换某一位后的通配键，distance 是转换层数。手算时按这个协议跟踪：hit 生成 *it、h*t、hi*，其中 h*t 桶找到 hot；hot 再通过 *ot 找到 dot 和 lot。

\textbf{代码落点。} 关键是访问标记在入队时设置。桶用完后可以清空，避免不同节点重复展开同一批邻居。关键代码行必须能逐句翻译成“旧状态怎样变成新状态”，否则就只是模板。

\textbf{正确性。} 每条转换边权都为 1，BFS 按层扩展；第一次到达 endWord 时经过的边数最少。

\textbf{边界实验。} 先用终点不在词表、起点和终点相同、桶重复扫描、单词长度一致做反例检查。实验时覆盖终点不在词表、起点和终点相同、桶重复扫描、单词长度一致，并打印入队时的完整状态。若同一位置但资源不同的状态被合并，很多图题会出现隐蔽错误。

\textbf{复杂度变化。} 两两建图 O(\ensuremath{N^{2}}L)，通配桶建图和 BFS 约 O(\ensuremath{NL^{2}}) 或按桶规模计，空间 O(NL)。

\textbf{迁移追问。} 如果题目改成“双向 BFS 可以进一步降低搜索空间”，先判断原状态是否仍然覆盖新答案集合，再决定是微调边界、改 value 语义，还是换算法。

\subsection{LeetCode 295 数据流的中位数}

\textbf{题目说明。} LeetCode 295 数据流的中位数 的题面入口要先还原成具体任务：数据流需要动态维护较小一半和较大一半。

\textbf{样例入口。} 拿数据流 1、2、3 手算，插入 1 后中位数 1；插入 2 后两半是 [1] 和 [2]，中位数 1.5；插入 3 后右半多一个，要平衡成左半 [2, 1]、右半 [3]。

\textbf{暴力基线。} 暴力每插入一个数就重新排序所有已到达元素，然后按长度奇偶取中位数。

\textbf{暴力过程。} 数据流 1、2、3 中，插入 1 中位数是 1；插入 2 后中位数是 (1+2)/2；插入 3 后中位数是 2。

\textbf{重复工作。} 题面模型：数据流需要动态维护较小一半和较大一半。暴力版的慢点要落到本题：浪费在每次都维护完整排序。中位数只需要较小一半的最大值和较大一半的最小值。后面的优化只消掉这类重复工作，不能改变答案集合。

\textbf{优化条件。} 本题的关键观察是：大顶堆保存较小一半，小顶堆保存较大一半，大小差不超过一。这句话必须能翻译成可维护状态；如果题目条件改变后观察不再成立，就不能继续套原来的写法。

\textbf{相邻状态。} 规律是大顶堆保存较小一半，小顶堆保存较大一半，大小差最多 1。把这个特例继续拆开看：堆顶必须正好代表下一步最需要处理的候选；若堆里可能有过期对象，读取堆顶前要先清理。堆只解决动态极值，不自动证明被弹出或被丢弃的元素无用。

\textbf{状态复用。} 大顶堆 small 保存较小一半，小顶堆 large 保存较大一半。保持 small.size() 与 large.size() 差不超过 1，且 small.top()<=large.top()。手算时只改被新输入影响的那一小部分，而不是回头重算完整候选。

\textbf{指针和字段。} 状态字段要能说清：small.top 是左半最大值，large.top 是右半最小值；奇数时较大堆顶是中位数，偶数时两个堆顶平均。手算时按这个协议跟踪：插入 1 放 small；插入 2 后移动平衡成 small=[1]、large=[2]；插入 3 后 large 多，再移动 2 到 small，中位数 2。

\textbf{代码落点。} 关键是先按顺序不变量放入堆，再重平衡大小。求平均时用 double，并防止两个 int 相加溢出。关键代码行必须能逐句翻译成“旧状态怎样变成新状态”，否则就只是模板。

\textbf{正确性。} 两个堆按大小分割数据流，中位数只可能位于左半最大或右半最小；大小不变量保证取堆顶即可。

\textbf{边界实验。} 先用偶数个元素、负数、重复值、平均值溢出做反例检查。实验时覆盖偶数个元素、负数、重复值、平均值溢出，记录堆顶是否过期、比较器是否让正确候选在堆顶、K 或窗口大小为边界值时是否稳定。

\textbf{复杂度变化。} 每次重新排序 O(n log n)；双堆 addNum O(log n)，findMedian O(1)，空间 O(n)。

\textbf{迁移追问。} 如果题目改成“若还要删除任意元素，需要延迟删除哈希表”，先判断原状态是否仍然覆盖新答案集合，再决定是微调边界、改 value 语义，还是换算法。

\subsection{LeetCode 303 区域和检索 - 数组不可变}

\textbf{题目说明。} LeetCode 303 区域和检索 - 数组不可变 的题面入口要先还原成具体任务：数组不变且要多次区间求和，单次扫描区间会重复读取同一段。

\textbf{样例入口。} 拿 nums=[-2, 0, 3, -5, 2, -1] 手算，sumRange(0, 2)=1，sumRange(2, 5)=-1。若每次查询都扫描区间会重复读同一段；预处理 prefix 后，闭区间和等于 prefix[right+1]-prefix[left]。

\textbf{暴力基线。} 慢版本按题意两两比较、枚举区间、扫描历史或持续迭代并查重。把检查条件写出来后，通常能看到当前状态只缺一个历史状态、规范键或已经访问过的中间状态；这正是从平方枚举或重复迭代走向哈希表的入口。放到本题就是：先按“数组不变且要多次区间求和，单次扫描区间会重复读取同一段”完整试慢版本；样例里已经能看到“规律是数组不可变时把所有前缀和提前缓存”，所以优化前必须先能说清慢版本枚举了哪些候选。

\textbf{暴力过程。} 拿 nums=[-2, 0, 3, -5, 2, -1] 手算，sumRange(0, 2)=1，sumRange(2, 5)=-1。若每次查询都扫描区间会重复读同一段；预处理 prefix 后，闭区间和等于 prefix[right+1]-prefix[left]。

\textbf{重复工作。} 题面模型：数组不变且要多次区间求和，单次扫描区间会重复读取同一段。暴力版的慢点要落到本题：慢版本会为“数组不变且要多次区间求和，单次扫描区间会重复读取同一段”反复寻找历史状态；而样例规律已经指出“规律是数组不可变时把所有前缀和提前缓存”。这些补数、前缀状态、规范键、半边和或访问过的中间状态可以用一个 key 归类，不必每次重新扫描或重复比较。后面的优化只消掉这类重复工作，不能改变答案集合。

\textbf{优化条件。} 本题的关键观察是：预处理 prefix[i+1] 为前 i+1 个数的和，sumRange(left, right) 用两个前缀相减。这句话必须能翻译成可维护状态；如果题目条件改变后观察不再成立，就不能继续套原来的写法。

\textbf{相邻状态。} 规律是数组不可变时把所有前缀和提前缓存。把这个特例继续拆开看：每个合法答案都要还原成当前状态和某个历史状态的配对；哈希表的 key 负责定位历史状态，value 负责表达次数、最早位置或存在性。先查还是先写，必须由是否允许当前前缀和自己配对来决定。

\textbf{状态复用。} 把回头扫描、两两比较或重复状态检测改成查表、分桶或访问记录；key 是当前题目需要的补数、前缀状态、规范表示、半边和或中间状态，value 根据目标选择次数、下标、列表或存在性。本题落点是：规律是数组不可变时把所有前缀和提前缓存。手算时只改被新输入影响的那一小部分，而不是回头重算完整候选。

\textbf{指针和字段。} 状态字段要能说清：当前状态或规范键、需要查询的历史 key、哈希表 value 语义、答案累计值或分桶结果；空前缀、空字符串、重复状态或初始状态要单独说明。手算时按这个协议跟踪：拿 nums=[-2, 0, 3, -5, 2, -1] 手算，sumRange(0, 2)=1，sumRange(2, 5)=-1。若每次查询都扫描区间会重复读同一段；预处理 prefix 后，闭区间和等于 prefix[right+1]-prefix[left]。然后按协议复查：手算时列出当前输入、当前状态或规范键、需要查询的历史 key、查到的 value、更新后的哈希表。这样能看出先查后插、先分桶后输出，还是先检测重复状态。

\textbf{代码落点。} 先查询还是先插入要由题意决定；key 构造必须无歧义；哈希表 value 的默认插入副作用要可控。关键代码行必须能逐句翻译成“旧状态怎样变成新状态”，否则就只是模板。

\textbf{正确性。} 证明从状态等价开始：任意合法答案都对应当前状态和某个历史 key、规范键或访问状态，查到的 key 也能反推出合法答案或合法分组。

\textbf{边界实验。} 先用left 为零、负数、单元素区间、多次查询、总和溢出做反例检查。实验时把left 为零、负数、单元素区间、多次查询、总和溢出加入对拍集合，用平方暴力和优化版比较。失败时先看初始历史状态、查询更新顺序和哈希表 value 类型。

\textbf{复杂度变化。} 暴力两两比较、枚举区间、扫描历史或重复迭代常见为平方级或更多；把历史状态、规范键或访问状态放进哈希表后，每步只做常数次查询和更新，通常是线性时间、线性空间。

\textbf{迁移追问。} 如果题目改成“若数组会更新，普通前缀和失效，需要树状数组或线段树”，先判断原状态是否仍然覆盖新答案集合，再决定是微调边界、改 value 语义，还是换算法。

\subsection{LeetCode 525 连续数组}

\textbf{题目说明。} LeetCode 525 连续数组 的题面入口要先还原成具体任务：把零看成负一后，零和一数量相等等价于前缀和相同。

\textbf{样例入口。} 拿 [0, 1, 0] 手算，把 0 当 -1，前缀差值序列 -1、0、-1；差值再次出现说明中间 0 和 1 数量相同。

\textbf{暴力基线。} 慢版本按题意两两比较、枚举区间、扫描历史或持续迭代并查重。把检查条件写出来后，通常能看到当前状态只缺一个历史状态、规范键或已经访问过的中间状态；这正是从平方枚举或重复迭代走向哈希表的入口。放到本题就是：先按“把零看成负一后，零和一数量相等等价于前缀和相同”完整试慢版本；样例里已经能看到“规律是保存差值第一次出现的位置，最长区间用最早位置”，所以优化前必须先能说清慢版本枚举了哪些候选。

\textbf{暴力过程。} 拿 [0, 1, 0] 手算，把 0 当 -1，前缀差值序列 -1、0、-1；差值再次出现说明中间 0 和 1 数量相同。

\textbf{重复工作。} 题面模型：把零看成负一后，零和一数量相等等价于前缀和相同。暴力版的慢点要落到本题：慢版本会为“把零看成负一后，零和一数量相等等价于前缀和相同”反复寻找历史状态；而样例规律已经指出“规律是保存差值第一次出现的位置，最长区间用最早位置”。这些补数、前缀状态、规范键、半边和或访问过的中间状态可以用一个 key 归类，不必每次重新扫描或重复比较。后面的优化只消掉这类重复工作，不能改变答案集合。

\textbf{优化条件。} 本题的关键观察是：哈希表保存每个前缀和最早下标，重复出现时更新最长长度。这句话必须能翻译成可维护状态；如果题目条件改变后观察不再成立，就不能继续套原来的写法。

\textbf{相邻状态。} 规律是保存差值第一次出现的位置，最长区间用最早位置。把这个特例继续拆开看：每个合法答案都要还原成当前状态和某个历史状态的配对；哈希表的 key 负责定位历史状态，value 负责表达次数、最早位置或存在性。先查还是先写，必须由是否允许当前前缀和自己配对来决定。

\textbf{状态复用。} 把回头扫描、两两比较或重复状态检测改成查表、分桶或访问记录；key 是当前题目需要的补数、前缀状态、规范表示、半边和或中间状态，value 根据目标选择次数、下标、列表或存在性。本题落点是：规律是保存差值第一次出现的位置，最长区间用最早位置。手算时只改被新输入影响的那一小部分，而不是回头重算完整候选。

\textbf{指针和字段。} 状态字段要能说清：当前状态或规范键、需要查询的历史 key、哈希表 value 语义、答案累计值或分桶结果；空前缀、空字符串、重复状态或初始状态要单独说明。手算时按这个协议跟踪：拿 [0, 1, 0] 手算，把 0 当 -1，前缀差值序列 -1、0、-1；差值再次出现说明中间 0 和 1 数量相同。然后按协议复查：手算时列出当前输入、当前状态或规范键、需要查询的历史 key、查到的 value、更新后的哈希表。这样能看出先查后插、先分桶后输出，还是先检测重复状态。

\textbf{代码落点。} 先查询还是先插入要由题意决定；key 构造必须无歧义；哈希表 value 的默认插入副作用要可控。关键代码行必须能逐句翻译成“旧状态怎样变成新状态”，否则就只是模板。

\textbf{正确性。} 证明从状态等价开始：任意合法答案都对应当前状态和某个历史 key、规范键或访问状态，查到的 key 也能反推出合法答案或合法分组。

\textbf{边界实验。} 先用全零、全一、从下标零开始、不要覆盖最早位置做反例检查。实验时把全零、全一、从下标零开始、不要覆盖最早位置加入对拍集合，用平方暴力和优化版比较。失败时先看初始历史状态、查询更新顺序和哈希表 value 类型。

\textbf{复杂度变化。} 暴力两两比较、枚举区间、扫描历史或重复迭代常见为平方级或更多；把历史状态、规范键或访问状态放进哈希表后，每步只做常数次查询和更新，通常是线性时间、线性空间。

\textbf{迁移追问。} 如果题目改成“这是前缀和最早位置语义，不是频次语义”，先判断原状态是否仍然覆盖新答案集合，再决定是微调边界、改 value 语义，还是换算法。

\begin{principlebox}{本节复盘方式}
不要只看最终算法名。每道题复盘时先遮住优化版，口述暴力枚举对象；再指出相邻窗口、历史前缀、候选区间、搜索状态或子问题里哪一部分被重复处理；最后用一个边界样例验证状态更新顺序。能讲完这三步，才算真正掌握本章案例。
\end{principlebox}
